%------------------------------------------------------------------------------%
% ALGEBRAIC TOPOLOGY                                                           %
%------------------------------------------------------------------------------%
% File  : Algebraic_Topology.tex                                               %
% Author: Klaus Hoeflinger, Beethovenstrasse 11, D-73117 Wangen                %
% Email : khoeflinger@t-online.de                                              %
%------------------------------------------------------------------------------%

%------------------------------------------------------------------------------%
%                       DOCUMENT CLASS and INCLUDE FILES                       %
%------------------------------------------------------------------------------%
\documentclass[11pt,twoside]{article}
\usepackage{geometry}
\geometry{paperheight=241mm,
          paperwidth=152mm,
          top=22mm,
          bottom=17mm,
          left=20mm,
          right=20mm,
          textwidth=112mm,
          textheight=202mm,
          headheight=10mm,
          headsep=3mm,
          footskip=9mm,
          includehead,
          includefoot,
          marginparsep=0mm,
          marginparwidth=0mm}

\renewcommand{\baselinestretch}{0.945}\normalsize

\AtBeginDocument{\setlength\abovedisplayskip{2mm}}
\AtBeginDocument{\setlength\belowdisplayskip{2mm}}

%------------------------------------------------------------------------------%
% define title formats                                                         %
%------------------------------------------------------------------------------%
\usepackage{titlesec}

\renewcommand{\thesection}{\arabic{section}}
\renewcommand{\sfdefault}{FiraSans}
\renewcommand{\ttdefault}{pcr}
\renewcommand{\rmdefault}{antiqua}

\titleformat{\part}[display]
{\normalfont \large \rmfamily \bfseries \centering}
{\small{\thepart.}}{2pt}{}

\titleformat{\section}[runin]
{\normalfont \rmfamily \bfseries}
{\thesection}{1pt}{.\;}[.]

%\titleformat{\section}
%{\normalfont \bfseries \small \filcenter}
%{\thesection.\;}{0mm}{}

\titleformat{\subsection}[runin]
{\normalfont \itshape}
{}{0mm}{}

\titleformat{\subsubsection}[runin]
{\normalfont \itshape}
{}{}{}{}

\titleformat{\paragraph}[runin]
{\normalfont}
{}{}{}{}

\titleformat{\subparagraph}[runin]
{\normalfont}
{}{}{}{}

\titlespacing*{\part}          { 0pt}{ 0pt}{ 8pt}
\titlespacing*{\section}       { 0pt}{ 7pt}{ 4pt}
\titlespacing*{\subsection}    { 0pt}{14pt}{ 6pt}
\titlespacing*{\subsubsection} { 0pt}{ 6pt}{12pt}
\titlespacing*{\paragraph}     { 0pt}{ 6pt}{12pt}
\titlespacing*{\subparagraph}  { 0pt}{ 0pt}{12pt}

%------------------------------------------------------------------------------%
% define enumerate parameters                                                  %
%------------------------------------------------------------------------------%
\usepackage{enumitem}
\setlength{\parindent}{3pt}
\setlength{\parskip}  {0pt}

%\setlist{noitemsep}

\setlist[enumerate]{
         leftmargin=12pt,
         rightmargin=0pt,
         itemsep=1pt,
         itemindent=3pt,
         topsep=3pt,
         labelsep=4pt,
         partopsep=0pt,
         labelwidth=0pt,
         listparindent=0pt,
         parsep=0pt}

\setlist[itemize]{
         leftmargin=12pt,
         rightmargin=0pt,
         itemsep=0pt,
         itemindent=1pt,
         topsep=1pt,
         labelsep=4pt,
         partopsep=1pt,
         labelwidth=0pt,
         listparindent=0pt,
         parsep=0pt}

%------------------------------------------------------------------------------%
% define packages                                                              %
%------------------------------------------------------------------------------%
\usepackage{upref}
\usepackage{amssymb}
\usepackage{slantsc}
\usepackage{amsmath}
\usepackage{amsthm}
\usepackage{thmtools}
\usepackage{amscd}
\usepackage{verbatim}
\usepackage{latexsym}
\usepackage{multicol}
\usepackage{graphicx}
\usepackage{setspace}
\usepackage[ngerman,english]{babel}
\usepackage[protrusion=true,expansion=true]{microtype}
\usepackage{tikz}
\usetikzlibrary{arrows,arrows.meta,decorations.markings}
\usepackage{mathrsfs}
\usepackage{amsfonts}
\usepackage{datetime2}
\usepackage{textgreek}
\usepackage{hyperref}
\usepackage{pifont}
\usepackage{relsize}
%\usepackage{biblatex}
\relsize{-0.05}
\usepackage[skip=0mm]{parskip}
\usetikzlibrary{decorations.pathmorphing}

\usepackage{mlmodern}
\usepackage[T5,T1]{fontenc}
\usepackage{ragged2e}

%\usepackage{mathabx}
%\usepackage{times}

%\usepackage{fontspec}
%\setmainfont{Nimbus Roman No9 L}
%\usepackage[T1]{fontenc}

%------------------------------------------------------------------------------%
% define page layout                                                           %
%------------------------------------------------------------------------------%
\usepackage{fancyhdr}
\pagestyle{fancy}
\setlength{\headheight}{30.0pt}
\renewcommand{\headrulewidth}{0pt}

\AtBeginDocument{
  \addtolength\abovedisplayskip{-0.0\baselineskip}
  \addtolength\belowdisplayskip{-0.0\baselineskip}
}

\fancyhead{}
\fancyfoot{}
\addtolength{\headwidth}{0.02\marginparwidth}

\makeatletter
\let\ps@plain\ps@fancy
\makeatother

\renewcommand\sectionmark[1]
{
 \def\sectionname{#1}
 \markright{\thesection #1}
}

\makeatletter
\@addtoreset{section}{part}
\makeatother

%------------------------------------------------------------------------------%
% define footnote without number                                               %
%------------------------------------------------------------------------------%
\newcommand{\myfootnote}[1]{%
\renewcommand{\thefootnote}{}%
\footnotetext{#1}%
\renewcommand{\thefootnote}{\arabic{footnote}}%
}

%------------------------------------------------------------------------------%
% define numbering in equations                                                %
%------------------------------------------------------------------------------%
\numberwithin{equation}{section}

%------------------------------------------------------------------------------%
% define newtheorem                                                            %
%------------------------------------------------------------------------------%
\declaretheoremstyle[
headfont=\scshape, 
headindent = 0mm, %\parindent,
%postheadhook = {\hspace{0mm}\newline},
%postheadspace = \newline, 
spaceabove = 3mm, 
spacebelow = 3mm,
numberwithin=section,
name=Theorem]{khMATH}
\declaretheorem[style=khMATH]{theorem}

\declaretheoremstyle[
headfont=\bfseries, 
headindent = 0mm, %\parindent,
%postheadhook = {\hspace{0mm}\newline},
%postheadspace = \newline, 
spaceabove = 3mm, 
spacebelow = 3mm,
numberwithin=part,
name=Theorem]{khMATH1}
\declaretheorem[style=khMATH1]{theorem1}

\declaretheoremstyle[
headfont=\scshape, 
headindent = 0mm, %\parindent,
%postheadhook = {\hspace{0mm}\newline},
%postheadspace = \newline, 
spaceabove = 3mm, 
spacebelow = 3mm,
numberwithin=section,
name=Corollary]{khMATH1}
\declaretheorem[style=khMATH1]{corollary}

%            name         counter     label              numeration-mode
\theoremstyle{plain}
\newtheorem*{introduction}            {\sc Introduction}
\newtheorem {standard}                {}                 [section]
\newtheorem*{definition}              {\sc Definition}
\newtheorem{satz1}                    {\sc Satz}
\newtheorem{lemma}                    {\sc Lemma}
\newtheorem*{proposition}             {\sc Proposition}
%\newtheorem{corollary}               {\sc Corollary}
\newtheorem*{corollar}                {\sc Korollar}
\newtheorem*{remark}                  {\sc Remark}
\newtheorem*{remarks}                 {\sc Remarks}
\newtheorem*{example}                 {Example}
\newtheorem*{examples}                {Examples}
\newtheorem*{theo}                    {\sc Theorem}

%------------------------------------------------------------------------------%
% define tabular                                                               %
%------------------------------------------------------------------------------%
\usepackage{tabularx}
\newcolumntype{L}[1]{>{\raggedright\arraybackslash}p{#1}}
\newcolumntype{C}[1]{>{\centering\arraybackslash}  p{#1}} 
\newcolumntype{R}[1]{>{\raggedleft\arraybackslash} p{#1}} 

%------------------------------------------------------------------------------%
% define \sh, \zB                                                              %
%------------------------------------------------------------------------------%
\def\dsh{{d.\,h.}}
\def\ie{{i.\,e.}}
\def\zB{z.\,B.}
\renewcommand{\abstractname}{ }

%------------------------------------------------------------------------------%
% change default spacing for cases                                             %
%------------------------------------------------------------------------------%
\makeatletter
\renewcommand*\env@cases[1][1.6]{%
  \let\@ifnextchar\new@ifnextchar
  \left\lbrace
  \def\arraystretch{#1}%
  \array{@{}l@{\quad}l@{}}%
}
\makeatother

\newenvironment{rcases}
  {\left.\begin{aligned}}
  {\end{aligned}\right\rbrace}

%------------------------------------------------------------------------------%
% change default for \predisplaypenalty                                        %
%------------------------------------------------------------------------------%
\predisplaypenalty=990
\allowdisplaybreaks \numberwithin{equation}{section}

%------------------------------------------------------------------------------%
% general notations                                                            %
%------------------------------------------------------------------------------%
\def\*{\star}
\def\={\,=\,}
\def\+{\,+\,}
\def\xsubset{{\,\subset\,}}
\def\xtimes{{\,\times\,}}
\def\xeof{{\,\in\,}}
\def\eq{{\,=\,}}
\def\xto{{\,\to\,}}
\def\eqdef{\,:=\,}
\def\:{{\,:\,}}
\def\ele{{\,\in\,}}
\def\exists{\mbox{ exists }}
\def\and{\mbox{ and }}
\def\for{\mbox{ for }}
\def\fa{\mbox{ for all }}
\def\fe{\mbox{ for each }}
\def\qfa{\quad \mbox{ for all }}
\def\df{\,\dotfill}
\def\iff{if and only if }
\def\wrt{with respect to}

\makeatletter
\newcommand \Df {\leavevmode \cleaders \hb@xt@ .70em{\hss .\hss }\hfill \kern \z@}
\makeatother

\def\sql{\mathfrak{a}}               % sesquilinear form a
\def\DF{B}                           % Dirichlet form a
%------------------------------------------------------------------------------%
% number sets and special elements                                             %
%------------------------------------------------------------------------------%
\def\P{\mathbb{H}}                   % half plane of $\C$
\def\J{I}                            % interval I
\def\N{{\mathbb{N}}}                 % unsigned integers
\def\Q{{\mathbb{Q}}}                 % rational integers
\def\Z{{\mathbb{Z}}}                 % integers
\def\R{{\mathbb{R}}}                 % real numbers
\def\Rn{{\mathbb{R}}^n}              % real numbers in Rn
\def\Rd{{\mathbb{R}}^d}              % real numbers in Rn
\def\C{{\mathbb{C}}}                 % complex numbers
\def\F{{\mathbb{K}}}                 % arbitrary field
\def\Torus{{\mathbb{T}}}             % complex circle line
\def\T{\tau}                         % upper bound of interval (0,t)
\def\sign#1{\operatorname{sign}(#1)} % sign of a number

%------------------------------------------------------------------------------%
% set theory                                                                   %
%------------------------------------------------------------------------------%
\def\G{G}                            % domain $G$
\def\clG{\overline{\G}}              % closure of $G$
\def\bndG{\partial \G}               % boundary of $G$
\def\setA{\mathcal{A}}               % set A
\def\setB{\mathcal{B}}               % set B
\def\setC{\mathcal{C}}               % set C
\def\setM{M}                         % set M
\def\setN{N}                         % set N
\def\setS{\mathcal{S}}               % set S
\def\famF{\mathfrak{F}}              % family of sets
\def\indI{\mathcal{I}}               % index set I
\def\pot#1{\mathfrak{P}(#1)}         % power set

\def\relR{\mathbf{R}}                % relation R
\def\mapf{f}                         % mapping f
\def\mapg{g}                         % mapping g

\def\set#1#2{\{ \, #1 \,: \, #2 \, \}}
\def\tensor#1#2{#1 \otimes #2}
\def\Tens{\oplus}

\def\cal#1{{\mathcal{#1}}}
\def\aut#1{\operatorname{Aut}(#1)}

\def\cross#1#2{#1 {\,\times\,} #2}
\def\multicross#1#2{#1 \times  \ldots \times  #2}
\def\multitens#1#2 {#1 \otimes \ldots \otimes #2}

%------------------------------------------------------------------------------%
% group theory                                                                 %
%------------------------------------------------------------------------------%
\def\1{{\mathsf{1}}}                % unit element of a monoid
\def\0{{\mathsf{0}}}                % unit element of an Abelian monoid
\def\kernel#1{\mathfrak{N}(#1)}     % kernel of a homomorphism
\def\cokernel#1{\mathfrak{C}{(#1)}} % cokernel of a homomorphism
\def\Hom#1{\operatorname{Hom}(#1)}

%------------------------------------------------------------------------------%
% linear algebra                                                               %
%------------------------------------------------------------------------------%
\def\vecV{\mathit{V}}               % vector space V
\def\vecH{\mathit{H}}               % vector space H
\def\vecW{\mathit{W}}               % vector space W
\def\vecU{\mathit{U}}               % subspace U
\def\convC{\mathcal{C}}             % convex set C
\def\basH{\mathfrak{B}}             % Hamel base B

\def\LO#1{\mathscr{L}(#1) }         % algebra of linear transformations

\def\scalar#1#2{ \langle #1,#2 \rangle }
\def\trace#1{\operatorname{tr}(#1)}
\def\dim#1{{\operatorname{dim} \, #1 }}
\def\codim#1{{\operatorname{codim} \, #1 }}
\def\dim{\operatorname{dim}}
\def\codim{\operatorname{codim}}
\def\ind{\operatorname{ind}}

\def\genG{\cal{G}}

\def\algA{\cal{A}}
\def\algB{\cal{B}}
\def\algU{\cal{U}}

\def\idI{\cal{I}}

%------------------------------------------------------------------------------%
% topology                                                                     %
%------------------------------------------------------------------------------%
\def\compK{{K}}          % compact set C
\def\openO{\mathcal{O}}  % open set O
\def\openU{\mathcal{U}}  % open set U
\def\U{\mathcal{U}}      % neighbourhood U
\def\V{\mathcal{V}}      % neighbourhood V
\def\d{\mathit{d}}       % metric function d

\def\interior#1{\mathring{#1}}
\def\closure#1{{\overline{#1}}}

\def\topT{\mathcal{T}}   % topology T
\def\topX{X}             % topological space X
\def\topY{Y}             % topological space Y
\def\topZ{Z}             % topological space Z
\def\metrS{M}            % metric space M

\def\grpG{G}             % topological group G
\def\grpA{A}             % topological group A
\def\grpB{B}             % topological group B
\def\grpC{C}             % topological group C
\def\grpH{H}             % topological group H
\def\grpU{U}             % topological group U
\def\grpV{V}             % topological group V
\def\topG{G}             % topological group G
\def\topH{H}             % topological group H
\def\topU{U}             % topological subgroup U

\def\subS{\cal{S}}
\def\riR{R}              % ring R

%------------------------------------------------------------------------------%
% calculus                                                                     %
%------------------------------------------------------------------------------%
\def\supp#1{\operatorname{supp}(#1)}
\def\singsupp#1{\operatorname{sing \, supp}(#1)}

\def\re{\operatorname{Re}}
\def\im{\operatorname{Im}}
\def\grad{\operatorname{grad}}

%\def\abs#1 {\left|#1\right|}
\def\abs#1 {\lvert #1 \rvert}
%\def\abs#1 {\left|\,#1\,\right|}

%------------------------------------------------------------------------------%
% measure theory                                                               %
%------------------------------------------------------------------------------%
\def\salgA{\Sigma}               % sigma-algebra S
\def\salgB{\cal{B}}              % sigma-algebra B
\def\Borel#1{\cal{B}_{#1}}       % Borel-algebra 
\def\LB{\lambda}                 % Lebesgue-Borel measure

%------------------------------------------------------------------------------%
% function spaces                                                              %
%------------------------------------------------------------------------------%
\def\Lp{L^p(\G)}                 % spaces L_p
\def\Lq{L^q(\G)}                 % spaces L_q
\def\Lh{L^2(\G)}                 % spaces L_2
\def\Li{L^{\infty}(\G)}          % spaces L_infty
\def\Lc{L_{loc}^1(\G)}           % spaces L_1^{loc}
\def\Ck{C^k(\G)}                 % spaces C_k

\def\BV#1{BV(#1) }               % set of functions of bounded variation
\def\AC#1{AC(#1) }               % set of absolutely continuous functions

\def\MP#1{\mathscr{M}^+(#1) }    % set of positive measures
\def\MS#1{\mathscr{M}^-(#1) }    % vector space of finite signed measures
\def\MC#1{\mathscr{M}(#1) }      % vector space of complex measures
\def\MCR#1{\mathscr{M}_{r}(#1) } % vector space of regular complex measures

\def\SobW{W^{k,p}(\G)}                % Sobolev space W
\def\SobO{\mathring{W}^{k,p}(\G)}     % Sobolev space W
%\def\WN#1#2{\mathring{W}^{#1}(#2) }   % Sobolev space W
\def\WN#1#2{W_0^{#1}(#2) }   % Sobolev space W
%\def\HN#1#2{\mathring{H}^{#1}(#2) }   % Sobolev space H
\def\HN#1#2{H_0^{#1}(#2) }   % Sobolev space H
%------------------------------------------------------------------------------%
% functional analysis                                                          %
%------------------------------------------------------------------------------%
\def\snorm#1{\lVert #1 \rVert}
\newcommand{\norm}[1]{\left\lVert #1 \right\rVert}

\def\topV{V}                       % topological vector space V
\def\topW{W}                       % topological vector space W
\def\normS{X}                      % normed space

\def\banX{\mathit{X}}              % Banach space X
\def\banY{\mathit{Y}}              % Banach space Y
\def\banZ{\mathit{Z}}              % Banach space Z
\def\dbanX{\mathit{X^*}}           % dual space of Banach space X
\def\dbanY{\mathit{Y^*}}           % dual space of Banach space Y
\def\dbanZ{\mathit{Z^*}}           % dual space of Banach space Z

\def\banA{\mathit{A}}              % Banach algebra A
\def\banB{\mathit{B}}              % Banach algebra B
\def\banC{\mathit{C}}              % C*-algebra C


\def\domain#1{\mathfrak{D}(#1)}    % domain
\def\range#1{\mathfrak{R}(#1)}     % range

\def\isoJ{\mathfrak{J}}            % isometry J
\def\repA{{\cal{A}_{\sql}}}        % representation operator A
\def\linS{{S}}                     % linear operator S
\def\linG{{G}}                     % linear operator G
\def\linL{{L}}                     % linear operator L
\def\linT{{T}}                     % linear operator A
\def\linA{\mathit{A}}              % linear operator A
\def\linB{{B}}                     % linear operator B
\def\linM{{M}}                     % linear operator M
\def\linU{{U}}                     % linear operator U
\def\linI{{I}}                     % linear operator I
\def\A{\cal{A}}                    % linear differential operator A
\def\BDO{\cal{B}}                  % linear boundary differential operator B
\def\linU{{U}}                     % unitary operator U
\def\linP{{P}}                     % projection P
\def\compA{k}                      % compact operator k
\def\linFR{f}                      % finite rank operator f
\def\fredA{A}                      % Fredholm operator F

\def\HAM{\cal{H}}                  % Hamilton operator
\def\PA{\partial^{\alpha}}         % elementary differential operator
\def\DO{\cal{A}}                   % linear differential operator
\def\BC{\cal{B}}                   % boundary operator
\def\SLO{\cal{L}_{p,q}}            % Sturm-Liouville operator
\def\MO{\cal{L}_{M}}               % momentum operator

\def\HS#1{S(#1)           }        % algebra of Hilbert-Schmidt operators
\def\FR#1{F(#1)           }        % algebra of finite rank operators
\def\TC#1{N(#1)           }        % algebra of trace class operators
\def\BU#1{\mathscr{U}(#1) }        % algebra of unitary linear operators
\def\BO#1{\mathscr{L}(#1) }        % algebra of bounded linear operators
\def\CO#1{\mathscr{C}(#1) }        % algebra of compact linear operators
\def\CL#1{\mathscr{C}(#1) }        % algebra of compact linear operators
\def\FO#1{\Phi(#1) }               % set of Fredholm operators
\def\FK#1#2{\Phi_{#2}(#1) }        % set of Fredholm operators with index k

\def\hilH{\mathit{H}}              % Hilbert space H
\def\clU{\mathit{U}}               % closed subspace U

\def\orthO{\mathfrak{O}}           % orthonormal system O
\def\orthS{\mathfrak{S}}           % orthonormal system S
\def\basB{\mathfrak{B}}            % basis B of a Hilbert space

\def\GL#1 {\mathbf{GL}(#1)}        % linear group
\def\OG#1 {\mathbf{O}(#1)}         % linear group
\def\SOG#1 {\mathbf{SO}(#1)}       % linear group

%------------------------------------------------------------------------------%
% distribution theory                                                          %
%------------------------------------------------------------------------------%
\def\disF{F}                  % distribution F
\def\disG{G}                  % distribution G
\def\disH{H}                  % distribution H
\def\disU{U}                  % distribution U
\def\disT{T}                  % distribution T
\def\SF{C^{\infty}(\G)}       % space of smooth functions
\def\TF{C_c^{\infty}(\G)}     % space of compactly supported smooth functions
\def\TFRd{C_c^{\infty}(\R^d)} % space of test functions
\def\SRd{\mathscr{S}(\R^d)}   % Schwartz space
\def\SR1{\mathscr{S}(\R)}     % Schwartz space
\def\B1#1{C^1(#1) }           % space of continuously differentiable functions
\def\Bm#1{C^m(#1) }           % space of continuously differentiable functions
\def\Ck{C^k(\G)}
\def\TD{\mathscr{S}'(\R^d)}   % space of temperated distributions
\def\E{\cal{E}}               % space of distributions with compact support
\def\D{\mathscr{D}}           % D
\def\FT{\mathcal{F}}          % Fourier transformation F
\def\FT{\mathscr{F}}          % Fourier transformation F

%------------------------------------------------------------------------------%
% other definitions                                                            %
%------------------------------------------------------------------------------%
\def\Proof{\textsc{Proof}}
\def\FRECHET{Fr\'{e}chet}
\def\ARZELA{Arzel\'{a}}
\def\GARDING{G\r{a}rding}
\def\HOELDER{H\"older}
\def\SCHROEDINGER{Schr\"odinger}
\def\JOERGENS{J\"orgens}
\def\WUEST{W\"uest}
\def\KAEHLER{K\"aehler}
\def\HOEFLINGER{H\"oflinger}
\def\POINCARE{Poincar\'{e}}
\def\FA{Functional Analysis}

\def\Author   {K.\,{\HOEFLINGER}}
\def\AuthorUC {K.\,HOEFLINGER}
\def\Version  {1.0}
\def\Email    {khoeflinger@t-online.de}


\def\E{E}
\def\B{B}
\def\simS{S}
\def\simK{K}
\def\simL{L}
\def\Rand#1{\delta_#1}
\def\modM{M}

\titleformat{\part}%[display]
{\normalfont \rmfamily \bfseries \centering}
{{\thepart.}}{3pt}{}

\titleformat{\section}[runin]
{\normalfont \bfseries \filcenter}
{\thesection.\,}{0mm}{}[.]

\titlespacing*{\part}   {5pt}{12pt}{ 4pt}
\titlespacing*{\section}{0pt}{ 3pt}{ 3pt}

\def\Title{Algebraic Topology}

%------------------------------------------------------------------------------%
%                                BEGIN OF DOCUMENT                             %
%------------------------------------------------------------------------------%
\begin{document}
\thispagestyle{empty}

\centerline{\large{\textbf{\Title}}}
\vspace{4pt}
\centerline{\small{{\textsc{\Author}}}}
\vspace{10pt}

\fancyhead[C] {\small{\textsc{\Title}}}
\fancyhead[OL]{\small{\thepage}}
\fancyhead[ER]{\small{\thepage}}

\myfootnote
{
  Version {\Version} from \today;
  Email:  \textit{khoeflinger@\,t-online.de}
}

%\begin{comment}
%------------------------------------------------------------------------------%
% content                                                                      %
%------------------------------------------------------------------------------%

\fancyhead[C] {\small{\textsc{I. Categories and Morphisms}}}
%------------------------------------------------------------------------------%
\part{Categories and Morphisms}
%------------------------------------------------------------------------------%
In category theory versucht man, 
die in verschiedenen mathematischen Disziplinen auftretenden structures
unter einem vereinheitlichenden Begriffssystem zu erfassen.

\section{Categories}
%------------------------------------------------------------------------------%
A \textit{category} is essentially the collection of mathematical structures
of the same kind.

\begin{definition}
A \textit{category} $\cal{C}$ is a system of mathematical objects $M,N,\ldots$,
together with sets $\operatorname{Mor}(M,N)$
for each two objects $M,N$ of $\cal{C}$, 
whose elements are called the \textit{morphisms} from $M$ to $N$.

Zusätzlich muß gelten:
for je drei objects $M_1$,$M_2$,$M_3$ von $\cal{C}$ existiert
eine Zuordnung 
\[
\circ: \cross{\operatorname{Mor}(M_1,M_2)}
             {\operatorname{Mor}(M_2,M_3)}
       \to    \operatorname{Mor}(M_1,M_3),
\]
die folgende Eigenschaften besitzt:
\begin{itemize}
\item
for $f \xeof \operatorname{Mor}(M_1,M_2)$, 
    $g \xeof \operatorname{Mor}(M_2,M_3)$, 
    $h \xeof \operatorname{Mor}(M_1,M_3)$
gilt
\[
h \circ (g \circ f) \eq (h \circ g) \circ f \quad \mbox{(Assoziativität)}
\]
\item
es existiert to each Objekt $M$ a morphism
$\1_M \xeof \operatorname{Mor}(M,M)$ with 
\[
\1_M \circ f \eq f \mbox{ and } g \circ \1_M \eq g 
\]
for all weiteren objects $N$ and 
alle Morphismen $f \xeof \operatorname{Mor}(M,N)$ and 
                $g \xeof \operatorname{Mor}(N,M)$.
\end{itemize}
\end{definition}

Häufig are the morphisms of a category
gerade the mappings between zwei objectsn,
which respect the natural structure of the objects.
Wir führen hierzu folgende Beispiele an:

\begin{examples}
\begin{itemize}
\item
The category \textit{(gr)} der groups, 
bestehend aus den groups and their homomorphisms.
\item
The category \textit{(ab)} der Abelschen groups, 
bestehend aus den Abelschen groups and their homomorphisms.
\item
The category \textit{(ri)} der Ringe and their homomorphisms.
\item
The category \textit{$Mod_R$} der $R$-modules over a ring $R$.
Ihre objects are the $R$-modules,
and the morphisms are the module-homomorphisms between ihnen.
\item
The category \textit{(vec)} der Vektorräume over a field $\F$.
Ihre objects are the Vektorräume selbst,
and the Morphismen are the $\F$-linear mappings.
Ebenso bilden the algebraischen Dualräume der Vektorräume über $\F$
eine category.
\item
The category \textit{($Alg_R$)} of $R$-algebras and their homomorphisms.
\item
The category \textit{(kk)} of chain complexes von $R$-modules and their 
homomorphisms.
\item
The topological category \textit{(top)}.
Ihre objects are the topological Räume,
ihre Morphismen are the continuous mappings between topological Räumen.
\item
The category der lokalkompakten topological groups,
whose morphisms are the continuous mappings as well.
\item
The differentialtopological category \textit{(difftop)},
bestehend aus den differenzierbaren Mannigfaltigkeiten
and den differenzierbaren mappings between ihnen.
\item
Es gibt categories,
deren Morphismen \textit{keine} mappings between ihren objectsn sind,
z.B. the homotopy category \textit{(htop)}. 
Ihre objects are the topological Räume,
ihre Morphismen jedoch the homotopy-Klassen $[X,Y]$ der stetiger mappings
$f:X \to Y$.
\end{itemize}
\end{examples}

{\bf Definition 2} (inverser morphism, Isomorphismus).

Ein morphism $g \xeof \operatorname{Mor}(N,M)$ heißt zum morphism 
$f \xeof \operatorname{Mor}(M,N)$ \textit{invers} oder a \textit{inverser morphism},
if $f \circ g \eq 1_N$ bzw. $g \circ f \eq 1_M$.
Those morphisms of a category,
which own an inverse morphism, 
are called the \textit{isomorphisms} of the category,
and in this case the related objects are called \textit{isomorph} to each another.

\begin{examples}
\begin{itemize}
\item
In der topological category \textit{(top)}
sind zwei topological spaces isomorphic,
wenn sie homeomorphic zueinander sind. 
\item
In der differentialtopological category \textit{(difftop)}
sind zwei Mannigfaltigkeiten isomorph,
wenn sie zueinander diffeomorph sind. 
\item
In der homotopy category \textit{(htop)} are zwei topological Räume
zueinander isomorph,
wenn sie homotopie-äquivalent sind.
\end{itemize}
\end{examples}

{\bf Definition 3} (Unterkategorie).

Let $\cal{C}$ a category.
Then heißt $\cal{C}$ a \textit{Unterkategorie}
einer weiteren category $\cal{D}$,
if each Objekt von $\cal{C}$ a Objekt von $\cal{D}$ ist
sowie each morphism $g \xeof \operatorname{Mor}(M,N)$ for objects $M,N$
von $\cal{C}$ a morphism der objects $M,N$ in $\cal{D}$ ist.

{\bf Definition 4} (volle Unterkategorie).

A Unterkategorie $\cal{C}$ von $\cal{D}$ heißt \textit{voll},
if each morphism von $\cal{D}$ auch a morphism von $\cal{C}$ ist.

\begin{example}
The category \textit{(ab)} der Abelschen groups is a volle Unterkategorie
der category \textit{(gr)} der groups.
\end{example}

\section{Functors}
%------------------------------------------------------------------------------%
\textit{functors} are mappings between den objectsn bzw. den Morphismen 
zweier categories,
die gewissen Eigenschaften genügen.

{\bf Definition 5} (functor).

Seien $\cal{C}$ and $\cal{D}$ zwei categories.
Ein \textit{kovarianter functor} $F$ between $\cal{C}$ and $\cal{D}$ 
ist a mapping, the 
\begin{itemize}
\item
each Objekt $M$ von $\cal{C}$ a Objekt $F(M)$ von $\cal{D}$ zuordnet
\item
each $f \xeof \operatorname{Mor}(M_1,M_2)$ a further morphism 
$f' \xeof \operatorname{Mor}(F(M_1),\F(M_2))$ zuordnet
such that \\
(a) $F(\1_M) \eq \1_{F(M)}$ \textit{(Identitätsaxiom)}\\
(b) $F(f \circ g) \eq F(f) \circ F(g)$ \textit{(chainsregel)}.
\end{itemize}

Ein functor $F$ heißt \textit{kontravariant},
if $F$ each morphism $f \xeof \operatorname{Mor}(M_1,M_2)$
a morphism $f' \xeof \operatorname{Mor}(F(M_2),F(M_1))$ zuordnet.

\begin{examples}
\begin{itemize}
\item
Let $\F$ a Körper and \textit{(vec)} the category der $\F$-Vektorräume
with den linear Operatoren als Morphismen and $(vec^{\*})$ the category
der Dualräume von $\F$-Vektorräumen.
Let $F$ be that mapping,
which assigns to each vector space $V$ its dual space $V^\*$.
Then each homomorphism $\phi:V \to W$ gives rise to a
homomorphism $\phi^\*:W^\* \to V^\*$ according to
$\phi^\* (w^\*) := v^\*$,
wobei $v^\*(v) := w^\*(\phi(v))$ for all $v \xeof V$.
Wir definieren $F(\phi) := \phi^\*$.
Then is $F$ a kontravarianter functor between den categories \textit{(vec)}
and $(vec^{\*})$.
\item
The \textit{Identitäts-functor}, der each Objekt and each morphism 
sich selbst zuweist.
\item
The \textit{Vergiß-functor} between den categories \textit{(ab)} and \textit{(gr)},
der each Abelian group sich selbst zuweist,
jedoch als Objekt in der category \textit{(gr)} der (allgemeinen) groups.
\end{itemize}
\end{examples}

\section{Natural Transformations}
%------------------------------------------------------------------------------%
Let $\cal{C}$ a category and seine $A,B$ zwei objects von $\cal{C}$.
Ist $\cal{D}$ a weitere category and are $F,G$ functors
von $\cal{C}$ nach $\cal{D}$,
so bilden these the objects $A,B$ in the objects $F(A),F(B)$
bzw. $G(A),G(B)$ der category $\cal{D}$ ab.
Ferner werden each morphism $\phi$ aus $\operatorname{Mor}(A,B)$
die Morphismen $F(\phi)$ aus $\operatorname{Mor}(F(A),F(B))$ bzw.
$G(\phi)$ aus $\operatorname{Mor}(G(A),G(B))$ zugeordnet.

{\bf Definition 6} (natural Transformation, natural Äquivalenz).

A \textit{natural Transformation} between $F$ and $G$ is a Familie
von Morphismen $\eta_A \xeof Mor(F(A),G(A))$ for each $A$ aus $\cal{C}$
such that
for all Morphismen $\phi \xeof Mor(A,B)$ gilt:
\[
G(\phi) \circ \eta_A \eq \eta_B \circ F(\phi).
\]
The functors $F$, $G$ are called zueinander \textit{natürlich äquivalent},
if a natural Transformation between ihnen existiert
and each morphism $\eta_A$ sogar a Isomorphismus ist.

If there is a natural transformation between the functors $F$ and $G$,
then the following diagram is commutative for each morphism $\phi \xeof Mor(A,B)$:

\begin{picture}(300,160)
\put( 90,15){$A$}
\put(250,15){$B$}
\put( 50,130){$F(A)$}
\put(110,90){$G(A)$}
\put(210,130){$F(B)$}
\put(270,90){$G(B)$}
\put(110, 18){\vector( 1, 0){130}}
\put( 80,133){\vector( 1, 0){120}}
\put(140, 93){\vector( 1, 0){120}}
\put( 90, 30){\vector(-1, 3){30}}
\put(250, 30){\vector(-1, 3){30}}
\put(100, 30){\vector( 1, 3){17}}
\put(260, 30){\vector( 1, 3){17}}
\put( 70,120){\vector( 2,-1){40}}
\put(230,120){\vector( 2,-1){40}}
\put(170,27){$\phi$}
\put(140,137){$F(\phi)$}
\put(190,97){$G(\phi)$}
\put( 95,113){$\eta_A$}
\put(255,113){$\eta_B$}
\end{picture}

\begin{example}
Let $\cal{C}$ the category \textit{(vec)} der $\R$-Vektorräume
and $\cal{D}$ the category der Bidualräume von $\R$-Vektorräumen.
Then are the functors $Id:V \mapsto V$ and $Bi:V \mapsto V^{\*\*}$
natürlich äquivalent.
The Morphismen $\eta_V$ are hierbei the Einbettungsabbildungen
$\iota:V \hookrightarrow V^{\*\*}$.
\end{example}

\fancyhead[C] {\small{\textsc{II. Homotopy and Fundamental Groups}}}
%------------------------------------------------------------------------------%
\part{Homotopy and Fundamental Groups}
%------------------------------------------------------------------------------%
By a homotopy is meant a continuous deformation from one mapping to another one.
any properties of mappings remain unchanged ander homotopies.
One can simplify the examination of a topological problem
by substituting a continuous mapping by a further but homotopic mapping.

\section{Homotopic Mappings}
%------------------------------------------------------------------------------%
Let us start with the precise definition of a homotopic mapping:

\begin{definition}
Let $\topX,\topY$ be topological spaces and $I := [0,1] \xsubset \R$.
Two mappings $f \:\topX \xto \topY$ and $g\:\topX \xto \topY$ are called
\textit{homotopic to each another} (notation: $f \simeq g$),
if there is a further continuous mapping $h \: \cross{I}{\topX} \xto \topY$
such that
$h(0,p) \eq f(p)$ and $h(1,p) \eq g(p)$ for all $p \xeof \topX$.
In this case $h$ is said to be a \textit{homotopy} between $f$ and $g$.
\end{definition}

\subparagraph{}
Thinking the parameter in the interval $[0,1]$ as the progress along the time,
then a homotopy $h$ between 
$f,g$ can be seen as follows:
at time $t \eq 0$ the mapping $h$ coincides with $f$
and is continuously deformated
such that in $t \eq 1$ it coincides with $g$.

\paragraph{}
Let us consider two simple examples:

\begin{itemize}[leftmargin=15pt]
\item[1.]
Especially each two continuous mappings $f_1,f_2 \: \Rd \xto \R$
are homotopic to each other,
since a homotopy $h$ is given by the mapping
$h \:(t,p) \mapsto t \cdot f_1(p) + (1-t) \cdot f_2(p)$.

\item[2.]
If a topological space $\topX$ consists of one element $p$ only,
then two mappings $f,g \: \topX \xto \topY$ are homotopic
iff
there is a continuous path $\gamma \:I \xto \topY$
with $\gamma(0) \eq f(p)$ and $\gamma(1) \eq g(p)$.
\end{itemize}

\subparagraph{}
The property of two mappings to be homotopic to each other
gives rise to an equivalence relation on the set of all continuous mappings
between topological spaces.
Let $[f]$ be the equivalence class of a continuous mapping $f\:\topX \xto \topY$,
then $[f]$ is called the \textit{homotopy class} of $f$.
It is usual to denote
as $[\topX,\topY]$ the totality of all homotopy classes of
$\topX$ and $\topY$\footnote
{
\,For punctured spaces and punctured continuous mappings
the homotopy classes $[(\topX,p),(\topY,q)]$ are defined similarly.}.

\subparagraph{}
For example,
the homotopy classes of the one-point space $\topX \eq \{p\}$
and a further topological space $\topY$ exactly correspond
to the set $\pi_0(\topY)$ of path-connected components of $\topY$.

\subparagraph{}
Forming the homotopy classes is compatible
with the composition of continuous mappings:
if $f,f'\:\topX \xto \topY$ and $g,g'\:\topY \xto \topZ$
are continuous with $f \simeq f'$ and $g \simeq g'$,
then the mappings $g \circ f$ and $g' \circ f'$ are homotopic to each other.
The topological spaces constitute a category
together with the homotopy classes $[\topX,\topY]$ as the morphisms.

\section{Homotopy Equivalences}
%------------------------------------------------------------------------------%
The notation of der \textit{homotopy equivalence} is of fundamental meaning
in algebraic topology,
sind it provides a way to classify topological spaces.
But this classification is not so exact as that given by homeomorphisms.
But both classifications are compatible in the sense
that homeomorphic spaces are also homotopy equivalent.

\subparagraph{}
Let $\topX,\topY$ be topological spaces and $Id_{\topX}, Id_{\topY}$
be the identity mappings.
\begin{itemize}[leftmargin=15pt]
\item[1.]
A continuous mapping $f:\topX \to \topY$ is called
a \textit{homotopy equivalence},
if there is a continuous mapping $g \:\topY \xto \topX$
such that
$g \circ f \simeq Id_{\topX}$ and $f \circ g \simeq Id_{\topY}$.
In this case $g$ is called a \textit{homotopy inverse} of $f$.
\item[2.]
The spaces $\topX,\topY$ are called \textit{homotopy-equivalent} rsp.
of one and the same \textit{homotopy type},
if there is a homotopy equivalence $f\:\topX \xto \topY$ 
(notation $\topX \simeq \topY$).
\end{itemize}

\subparagraph{}
If $g$ is a homotopy inverse of $f$,
then $f$ is a homotopy inverse of $g$,
and the mappings $f$ and $g$ are called 
\textit{inverse homotopy equivalences} in this case.
For example,
two homeomorphic spaces are homotopy-equivalent to each other.
But the reverse statement does not hold:
the sets $\{0\}$ and $\R$ are merely homotopy-equivalent,
but not homeomorphic.

\paragraph{}
A topological space $\topX$ is called \textit{contractible},
if it is homotopy-equivalent to the singleton $\{p\}$.
The following examples for contractible spaces are well-known:

\begin{itemize}[leftmargin=15pt]
\item[1.]
The Euclidean spaces $\Rd (d \xeof \N)$ are contractible.
\item[2.]
A subset $S \xsubset \Rd$ is called \textit{star-shaped},
if there is a point $p \xeof S$
such that
for each further point $q \xeof S$
the segment $\vec{pq}$ is also a subset of $S$.
Each star-shaped set is contractible.
\end{itemize}

\section{Retractions}
%------------------------------------------------------------------------------%
The homotopy-equivalence between a topological space
and one of its subspaces cannot always immediately be recognized.
Thus one is searching for criterions,
which ensures the homotopy-equivalence of a space $\topX$
with a subspace $\setM \subset \topX$.
One of them is the concept of \textit{retraction}.

\subparagraph{}
Let $\topX$ be a topological space and $\setM \xsubset \topX$ be a subset.
A \textit{retraction} of $\topX$ in $\setM$ is a continuous mapping
$\rho \: \topX \xto \setM$ with $\rho(p) \eq p$ for $p \xeof \setM$.
The set $\setM$ is called a \textit{retract} of $\topX$
if there is a retraction of $\topX$ in $\setM$.

xxx
\subparagraph{}
A Verst\"arkung des Rektrakt-Begriffes is derjenige
des \textit{Deformationsretraktes}.

{\bf Definition 4} (Deformationsretrakt).

Let $\topX$ a topological space and $\setM \subset \topX$.
A Retraktion $\rho:\topX \to \setM$ is called a \textit{Deformationsretraktion},
if $\rho$ homotop zur Identit\"atsabbildung $Id:\topX \to \topX$ ist.
The subset $\setM$ is called in diesem Falle a \textit{Deformationsretrakt}.

\subparagraph{}
A Deformationsretraktion $\rho:\topX \to \setM$ and the Inklusionsabbildung
$\iota:\setM \to \topX$ are also zueinander inverse homotopy-equivalenceen.
Zwei spaces are daher homotopie-\"aquivalent,
if der a space a Deformationsretraktion des anderen space ist.
Aus diesem Grand are Deformationsretraktionen von gro\"sser Bedeutung.
Sie are oft leichter zu erkennen als a Paar von
homotopie-\"aquivalenten mappings.

{\bf Definition 5} (starke Deformationsretraktion).

A Retraktion $\rho:\topX \to \setM$ is called a \textit{starke}
Deformationsretraktion,
if for the homotopy $h$ between $\rho$ and $Id_{\topX}$ gilt:
$h(t,p) \eq p$ for all $t$ and $p \xeof \setM$.

The subset $\setM$ is called in diesem Falle ein
\textit{starker Deformationsretrakt} von $\topX$.
Each point in $\setM$ bleibt hierbei w\"ahrend der Deformation fest.

\begin{example}
The Nullpunkt $\{0\} \xeof \R^n$ is a starker
Deformationsretrakt sowohl der Vollkugel $D^n$ als auch des space $\R^n$.
\end{example}

\section{The Fundamental Group of a Topological Space}
%------------------------------------------------------------------------------%
The \textit{Fundamentalgruppe} is das einfachste algebraische Objekt,
das man einem topological space zuweisen kann.
Ihre Definition beruht auf dem Begriff der \textit{nullhomotopen Schleifen}
an einem point.

Let $(\topX,p)$ a punktierter topological space.
Then is called each (geschlossene) continuous Weg $\gamma:S^1 \to \topX$ mit
$\gamma(1) \eq p$ a \textit{Schleife} an $p$.

A Schleife is called \textit{nullhomotop},
if sie homotop zum \textit{konstanten} Weg $\gamma_0$ with $\gamma_0(t) \eq p$
for all $t \xeof I$ ist.
Nullhomotope Schleifen lassen sich also stetig zu ihrem Grundpunkt
zusammenziehen.
Dawith hat man a Instrument zur Hand,
um \textit{L\"ocher} in einem topological space mathematisch exakt zu erfassen.

Let $(\topX,p)$ a punktierter space and $\Omega(\topX,p)$
die set der Schleifen an $p$.
Auf $\Omega(\topX,p)$ l\"a\"sst sich a innere Verkn\"upfung $\circ$ wie folgt
definieren:
for $\gamma_1, \gamma_2 \xeof \Omega(\topX,p)$ is the Schleife $\gamma_2 \circ
\gamma_1$ the Hintereinanderausf\"uhrung der beiden Schleifen.
The zu einer Schleife $\gamma$ inverse Schleife $\gamma^{-1}$ ist
diejenige Schleife, the man erh\"alt,
indem $\gamma$ \textit{r\"uckw\"arts} durchlaufen wird.

Let $\pi_1(\topX,p) := [(S^1,1),(\topX,p)]$ the set der homotopy classes
von Schleifen an $p$.
Then induziert the auf $\Omega(\topX,p)$ erkl\"arte Verkn\"upfung $\circ$ eine
Verkn\"upfung auf $\pi_1(\topX,p)$ according to
$[\gamma_1] \circ [\gamma_2] := [\gamma_1 \circ \gamma_2]$ for
$[\gamma_1], [\gamma_2] \xeof \pi_1(\topX,p)$.
Diese Definition is unabh\"angig von den gew\"ahlten Repr\"asentanten der
homotopy classes.

Ein neutrales Element $e \xeof \pi_1(\topX,p)$ is diejenige homotopyklasse,
die den konstanten Weg $\gamma_0$ with $\gamma_0(t) \eq p$ for $t \xeof I$
(Nullschleife) enth\"alt.
\"Uberdies gibt es zu each $[\gamma] \xeof \pi_1(\topX,p)$ a inverses Element
$[\gamma]^{-1}$
such that
$[\gamma] \circ [\gamma]^{-1} \eq e$.

Hence the set $\pi_1(\topX,p)$ tr\"agt the structure of a (in general)
non-commutative group.

{\bf Definition 6} (Fundamentalgruppe).

The group $(\pi_1(\topX,p),\circ)$ is called the \textit{Fundamentalgruppe} oder
\textit{erste homotopygruppe} des punktierten space $(\topX,p)$.

Beispiel: sei $(\topX,p) \eq (\R,0)$.
Then is each Schleife an $0$ nullhomotop.
Es gibt daher genau a homotopyklasse von Schleifen an $0$,
and hence we have $\pi_1(\R,0) \cong \{0\}$.

\section{The Fundamental Group Functor}
%------------------------------------------------------------------------------%
Ist $f:(\topX,p) \to (\topY,q)$ a punktierte continuous mapping and
$\gamma$ a Schleife an $p$,
so is the mapping $f \circ \gamma$ a Schleife an $q$.
Ferner werden zueinander homotope Schleifen an $p$ unter $f$ in
zueinander homotope Schleifen an $q$ \"ubergef\"uhrt.

The group operations of the fundamental group $\pi_1(\topX,p)$
and $\pi_1(\topY,q)$ are with $f$ insofern vertr\"aglich,
als dass $f(\gamma_2 \circ \gamma_1) \eq f(\gamma_2) \circ f(\gamma_1)$ gilt.
The mapping $f$ induziert daher a group homomorphism
$f_{\*}: \pi_1(\topX,p) \to \pi_1(\topY,q)$.

Hieraus folgt,
dass the Zuordnungen $(\topX,p) \mapsto \pi_1(\topX,p)$ and $f \mapsto f_{\*}$
einen kovarianten functor $\pi_1$ von der category der
punktierten topological spaces in the category der groups darstellen.
Dieser functor $\pi_1$ is called daher \textit{Fundamentalgruppen-Funktor}
bzw.  \textit{erster Homotopie-Funktor}.

A wichtige Eigenschaft zweier homotoper mappings ist,
dass sie identische group homomorphisms between den entsprechenden
Fundamentalgruppen erzeugen:

{\bf Satz 1}

Seien $f,g:\topX \to \topY$ zueinander homotope mappings.
Then gilt $f_{\*} \eq g_{\*}$.

Desweiteren gilt,
dass a homotopy-equivalence between wegzusammenh\"angenden
topological spaces die
isomorphy ihrer Fundamentalgruppen zur Folge hat.

{\bf Satz 2}

Seien $\topX,\topY$ zwei wegzusammenh\"angende topological spaces.
Sind $\topX$ and $\topY$ zueinander homotopie-\"aquivalent,
so are the Fundamentalgruppen $\pi_1(\topX)$ and $\pi_1(\topY)$
zueinander isomorph.

Da dies insbesondere for hom\"oomorphe spaces gilt,
ist the Fundamentalgruppe einer Klasse zueinander hom\"oomorpher spaces
bis auf isomorphy eindeutig bestimmt.

Let $\topX$ a topological space and $p,q \xeof \topX$ zwei points.
Es stellt sich sofort the Frage,
inwiefern the Fundamentalgruppen $\pi_1(\topX,p)$ and $\pi_1(\topX,q)$
vergleichbar sind.
Klar ist,
dass for points aus verschiedenen Zusammenhangskomponenten
des space keinerlei Beziehungen between ihren Fundamentalgruppen
existieren k\"onnen.

Ist $\topX$ jedoch \textit{wegzusammenh\"angend},
then gibt es einen group Isomorphismus between den Fundamentalgruppen
$\pi_1(\topX,p)$ and $\pi_1(\topX,q)$.
Jeder continuous Weg between $p$ and $q$ induziert einen solchen Isomorphismus,
ein \textit{kanonischer} Isomorphismus liegt daher nicht vor.
For einen wegzusammenh\"angenden space $\topX$ kann man also von \textit{der}
Fundamentalgruppe $\pi_1(\topX)$ des space sprechen.

A gewisse Sonderstellung nehmen diejenigen wegzusammenh\"angenden spaces ein,
deren Fundamentalgruppen nur aus dem neutralen Element bestehen.

{\bf Definition 7} (einfach zusammenh\"angender space).

Ein wegzusammenh\"angender topological space $\topX$ is called
\textit{einfach zusammenh\"angend},
if seine Fundamentalgruppe trivial ist,
that is,
$\pi_1(\topX) \cong \{0\}$.

The simpelsten Beispiele for einfach zusammenh\"angende spaces are the
kontrahierbaren spaces.

\section{The Computation of Fundamental Groups}
%------------------------------------------------------------------------------%
Fundamentalgruppen are im allgemeinen nicht einfach zu berechnen.
Ein erstes grundlegendes Resultat,
dass man am einfachsten with Methoden der Funktionentheorie beweist,
ist the Bestimmung der Fundamentalgruppe des Kreises:

{\bf Satz 3}

$\pi_1(S^1) \cong \Z$.

For a kartesisches Produkt von punktierten topological spaces
gewinnt man the Fundamentalgruppe durch direkte Produktbildung der
Fundamentalgruppen der Faktorr\"aume.

{\bf Satz 4}

Seien $(\topX,p),(\topY,q)$ punktierte topological spaces.
Then gilt
\[
\pi_1( \cross{\topX}{\topY}, (p,q))
\, \cong \,
\pi_1(\topX,p) \times \pi_1(\topY,q).
\]

\begin{example}
The Fundamentalgruppen der $n$-dimensionalen Tori sind
$\pi_1(\T^n) \eq \pi_1({S^1} \times \ldots \times {S^1}) \cong \Z^n$.
\end{example}

The folgende Satz,
ein Spezialfall des Satzes von \textit{Seifert-van Kampen},
gestattet the Berechnung der Fundamentalgruppe of a space,
sofern man ihn in geeignete Teilr\"aume zerlegen kann
and the Fundamentalgruppen dieser Teilr\"aume kennt.

{\bf Satz 5} (Variante des Satzes von Seifert-van Kampen).

Let $\topX$ a topological space and $A,B \subset \topX$
offene and zusammenh\"angende subsets
such that
$A \cup B \eq \topX$ and $A \cap B$ einfach zusammenh\"angend ist.
Then gilt for einen point $p \xeof A \cap B$
\[
\pi_1(\topX,p) \, \cong \, \pi_1(A,p) \* \pi_1(B,p).
\]

Hierbei is $\*$ das freie Produkt der groups $\pi_1(A,p)$ and $\pi_1(B,p)$.

\begin{example}
The spaces $S^n$ erf\"ullen the Voraussetzungen dieses Satzes.
Speziel for den Fall $n=2$ erh\"alt man a Zerlegung der Kugeloberfl\"ache $S^2$,
wenn man vom Nordpol and S\"udpol aus gen\"ugend gro\"sse Umgebungen definiert,
such that
whose intersection a einfach zusammenh\"angende Umgebung
des \"Aquators darstellt.
Es gilt daher $\pi_1(S^2) \cong \{0\}$.
Dieselbe Argumentation bleibt auch for $n \geq 3$ g\"ultig.
\end{example}

\section{Higher Homotopy Groups}
%------------------------------------------------------------------------------%
Ersetzt man den Definitionsbereich $S^1$ von Schleifen an einem point
$p \xeof \topX$ durch the h\"oherdimensionalen Sph\"aren $S^n$,
so erh\"alt \textit{h\"ohere homotopy groups} als Verallgemeinerung
des Fundamentalgruppen-Begriffes.

{\bf Definition 8} ($n$-te homotopy-group).

Let $N$ be the \textit{Nordpol} of the $n$-dimensional sphere $S^n$ $(n \geq 2)$,
$(\topX,p)$ a punctured topological space and
$[S^n,(\topX,p)]$ the set der homotopy classes of punctured
continuous mappings $f:(S^n,N) \to (\topX,p)$.
Then is called $\pi_n(\topX,p) := [S^n,(\topX,p)]$ the $n$-te homotopy group of
$(\topX,p)$.

For $n=1$ stimmt these Definition with derjenigen der Fundamentalgruppe
$\pi_1(\topX,p)$ \"uberein.
A group structure von $\pi_n(\topX,p)$ for $n \geq 2$ erh\"alt man wie folgt:

Da $S^n$ hom\"oomorph zum Quotientenraum $I^n / \partial I^n$ ist,
gibt es zu each continuous mapping $f:(S^n,N) \to (\topX,p)$
genau a continuous mapping $f': I^n \to \topX$ with ${f'}^{-1}(p) \eq \partial I^n$.
For $[f],[g] \xeof \pi_n(\topX,p)$ betrachte the zugeh\"origen mappings
$f',g':I^n \to \topX$.
Definiere $f' \* g':I^n \to \topX$ durch Halbieren and Zusammensetzen der
jeweiligen Definitionsbereiche.
Da $f'(x) \eq g'(x)$ for Randpunkte $x$ des Einheitsw\"urfels gilt,
ist $f' \* g'$ in points der Trennfl\"ache stetig.
Then is for $f$,$g$ the innere Verkn\"upfung $f \star g$ definiert according to
$f \star g := [f' \* g']$.
The Verkn\"upfung is unabh\"angig von den gew\"ahlten Repr\"asentanten.

H\"ohere homotopy-groups besitzen folgende Eigenschaft:

{\bf Satz 5}

Let $\topX$ a topological space and $n \geq 2$.
Then is $\pi_n(\topX)$ an Abelian group.

Let $\topX$ a wegzusammenh\"angender topological space and
$\topX^{\sim}$ seine topological \"Uberlagerung.
Then gilt $\pi_n(\topX) \eq \pi_n(\topX^{\sim})$ for $n \geq 2$.

Im Verallgemeinerung des einfach zusammenh\"angenden space is called ein
topological space $n$-fach zusammenh\"angend,
if $\pi_i(\topX) \cong 0$ for $i=0,\ldots,n$.

Speziell for the topological spaces $S^m$ erh\"alt man the homotopy groups
$\pi_n(S^m)$.
Ihre Berechnung is nicht trivial and bis heute noch nicht abgeschlossen.
The wichtigsten Resultate are in folgendem Satz zusammengefa\"sst.

{\bf Satz 6}

\begin{itemize}
\item
$\pi_n(S^n) \cong \Z$ for $n \xeof \N$.
\item
$\pi_n(S^m) \cong 0$ for $n > m$.
\item
$\pi_3(S^2) \cong \Z$. \textit{(Hopf-Faserung).}
\item
$\pi_1(S^m) \cong 0$ for $m \geq 2$,
that is,
the spaces $S^m$ are for $m \geq 2$ einfach zusammenh\"angend.
\end{itemize}

Speziell for $n \eq m$ gibt es a bijektive mapping between den
homotopy classes $[S^n,S^n]$ and $\Z$.
Dawith is each continuous mapping $f:S^n \to S^n$ a ganze Zahl zugeordnet,
welche der \textit{mappingsgrad} $\operatorname{grad}(f)$ von $f$ genannt wird.

\begin{example}
Let $A: S^n \to S^n$ the antipodische mapping with $A(x) := -x$.
Then besitzt $A$ den mappingsgrad $\operatorname{grad}(A) \eq (-1)^{n+1}$.
\end{example}

For zwei mappings $f,g:S^n \to S^n$ gilt:
$\operatorname{grad} (f \circ g) =
 \operatorname{grad} (f) \cdot
 \operatorname{grad} (g)$.
Ist $f$ nicht surjektiv, so is $\operatorname{grad} (f) \eq 0$.

\fancyhead[C] {\small{\textsc{III. Ueberlagerungen}}}
%------------------------------------------------------------------------------%
\part{Ueberlagerungen}
%------------------------------------------------------------------------------%

\section{Basic Concepts}
%------------------------------------------------------------------------------%
The Begriff der \textit{topological \"Uberlagerung} trat historisch zum ersten
al in der Funktionentheorie auf.

{\bf Definition 1} (topological \"Uberlagerung).

Seien $E,B$ topological spaces and $\pi:E \to B$ a surjektive stetige
mapping.
The space $E$ is called a \textit{(topologische) \"Uberlagerung} von $B$,
wenn es zu each point $q \xeof B$ a Umgebung $V_q$ gibt,
such that

\begin{itemize}
\item
$V_q$ is weg-zusammenh\"angend.
\item
das Urbild $\pi^{-1}(V_q)$ is the disjunkte union von h\"ochstens
abz\"ahlbar vielen offenen sets $\topU_{\alpha} \subset E$,
die zu $V_q$ jeweils hom\"oomorph are \textit{(lokale Trivialisierung)}.
\end{itemize}

Es gilt also in diesem Fall $\pi^{-1}(V_q) \cong V_q \times I$
with einer h\"ochstens abz\"ahlbaren
and with der diskreten Topologie versehenen Indexmenge $I$.
Aus dieser Definition folgt,
dass es zu each $p \xeof E$ a offene Umgebung $U_p \subset E$ gibt,
die hom\"oomorph zur set $\pi(U_p) \subset B$ ist.
In diesem Sinne is $E$ \textit{lokal hom\"oomorph} zu $B$.

The mapping $q \mapsto \operatorname{Card}(\pi^{-1}(q))$,
die each point $q \xeof B$ the M\"achtigkeit ihrer Faser zuordnet,
ist lokal konstant.
Ist a \"Uberlagerungsraum $E$ zusammenh\"angend,
so haben all Fasern $\pi^{-1}(q)$ ($q \xeof B$) dieselbe M\"achtigkeit.
The Kardinalzahl der Fasern $\pi^{-1}(q)$ is called in diesem Fall
die \textit{Bl\"atterzahl} der \"Uberlagerung.

\begin{examples}
\begin{itemize}
\item
Let $\topX$ a h\"ochstens abz\"ahlbare set,
versehen with der diskreten Topologie.
Then is $\cross{\topX}{B}$ with der Projektionsabbildung
eine covering von $B$.
\item
Let $E \eq \R$, $B \eq S^1 \subset \C$ and $\pi(x) := e^{ix}$, $x \xeof \R$.
Then stellt der space $\R$ a covering von $S^1$ dar.
The Urbilder von offenen sets in $S^1$ are abz\"ahlbare unions
von offenen Intervallen $(a,b) \subset \R$.
\item
Let $E \eq S^n$, $B \eq \R\,P^n$ and $\pi(p) := \{p,-p\}$.
Then is $(S^n,\pi,\R\,P^n)$ a zweibl\"attrige \"Uberlagerung
des projektiven space.
\end{itemize}
\end{examples}

\section{The Lifting of Paths}
%------------------------------------------------------------------------------%
Das \textit{Hochheben} von Wegen $\gamma:I \to B$ in den
\"Uberlagerungsraum $E$ is der Kern der \"Uberlagerungstheorie.

{\bf Definition 2} (Hochhebung).

Let $\pi:E \to B$ a covering with $\pi(p) \eq q$, $I := [0,1]$ and
$\gamma: I \to B$ a stetiger Weg with $\gamma(0) \eq q$.
A \textit{Hochhebung} von $\gamma$ is a stetiger Weg
$\gamma': I \to E$ with $\gamma'(0) \eq p$ and $\pi \circ \gamma' \eq \gamma$.

\begin{picture}(200,90)
\put(130,75){$(E,p)$}
\put( 60,10){$I$}
\put(130,10){$(B,q)$}
\put( 70,12){\vector( 1, 0){55}}
\put(140,70){\vector( 0,-1){45}}
\put( 70,15){\vector( 1, 1){57}}
\put(100,20){$\gamma$}
\put(145,45){$\pi$}
\put( 80,45){$\gamma'$}
\end{picture}

For Hochhebungen erh\"alt man den folgenden Existenz- and Eindeutigkeitssatz.

{\bf Satz 1}

Let $\pi:E \to B$ a covering with $\pi(p) \eq q$.
Zu each continuous Weg $\gamma: I \to B$ with $\gamma(0) \eq q$ gibt es genau
eine Hochhebung $\gamma'$ in den \"Uberlagerungsraum $E$ with $\gamma'(0)=p$.

Sind $\gamma$ and $\delta$ zwei homotope Wege in $B$ with gleichem
Anfangspunkt $q \xeof B$,
then are auch ihre Hochhebungen nach $E$ homotop.
The homotopy von Wegen bleibt also bei Hochhebungen erhalten.
Dies is the Aussage des \textit{Monodromie-Lemmas}.
Insbesondere besitzen the Hochhebungen von $\gamma$ and $\delta$
den identischen Endpunkt $p \xeof E$
such that
dieser durch the homotopy-Klasse $[\gamma]$ eindeutig bestimmt ist.

Dawith is a Operation der Fundamentalgruppe $\pi_1(B,q)$ auf der
Faser von $q$ erkl\"art.
Ist der \"Uberlagerungsraum $E$ wegzusammenh\"angend,
so operiert the group $\pi_1(B,q)$ sogar transitiv auf $\pi^{-1}(q)$.

The \"Uberlagerungstheorie von topological spaces is daher eng mit
der Theorie der Fundamentalgruppen verkn\"upft.

\section{Charakteristic Subgroups}
%------------------------------------------------------------------------------%
Covering mappings induce group homomorphisms between the
fundamental groups of the total space and the base space.

{\bf Satz 2}

Let $\pi:\E \to \B$ a covering with $\pi(p) \eq q$.
Then is der von $\pi$ induzierte group homomorphism
$\pi_{\*}: \pi_1(\E,p) \to \pi_1(\B,q)$ injektiv.

For a covering besitzt daher the Fundamentalgruppe des
\"Uberlagerungsraumes a einfachere structure than the base space.

{\bf Definition 3} (charakteristische Untergruppe).

Das Bild $G(\E,p,\pi)\,:=\, \pi_{\*}(\pi_1(\E,p)) \subset \pi_1(\B,q)$
is called the \textit{charakteristische Untergruppe} der \"Uberlagerung.

Zwischen der group-Operation von $\pi_1(\B,q)$ auf $\pi^{-1}(q)$ and der
charakteristischen subgroup $G(\E,p,\pi)$ besteht folgender Zusammenhang.

{\bf Satz 3}

Let $\pi:\E \to \B$ a covering with $\pi(p) \eq q$. Then gilt
\begin{itemize}
\item
The Isotropie-Gruppe von $p \xeof \pi^{-1}(q)$ is $G(\E,p,\pi)$.
\item
Ist $\E$ wegzusammenh\"angend, so are for $p,p' \xeof \pi^{-1}(q)$
die Untergruppen $G(\E,p,\pi)$ and $G(\E,p',\pi)$
als Untergruppen von $\pi_1(\B,q)$ zueinander konjugiert.
\end{itemize}

{\bf Definition 4} (\"aquivalente coverings).

Zwei coverings $\pi:\E \to \B$ and $\pi':\E' \to \B$
eines topological space $\B$ are called \textit{\"aquivalent},
if es einen Hom\"oomorphismus $h:\E \to \E'$ gibt
such that $\pi \eq \pi' \circ h$.

Dawith lassen sich coverings \"uber einem fest gew\"ahlten Basisraum
klassifizieren.

{\bf Satz 3}

Zwei wegzusammenh\"angende coverings von $\B$ are genau then \"aquivalent,
wenn ihre charakteristischen Untergruppen in $\pi_1(\B)$ zueinander
konjugiert sind.

\section{Universal Coverings}
%------------------------------------------------------------------------------%
\textit{Einfach zusammenh\"angende} \"Uberlagerungsr\"aume are the strukturell
einfachsten coverings, the es gibt.
Sie werden daher wie folgt ausgezeichnet:

{\bf Definition 5} (universal \"Uberlagerung).

A covering $\pi:\E \to \B$ is called \textit{universell},
if $\E$ einfach zusammenh\"angend ist.

Ist a \"uberlagerter space $\B$ lokal wegzusammenh\"angend,
so are zwei universal coverings $\E$ and $\E'$ zueinander \"aquivalent.
an kann in diesem Fall daher von \textit{der} universellen covering von $\B$
sprechen.

Es stellt sich the Frage,
unter welchen Voraussetzungen a universal covering of a lokal
wegzusammenh\"angenden topological space $\B$ existiert.

{\bf Definition 6} (semi-lokal einfach zusammenh\"angender space).

Ein topological space $\B$ is called \textit{semi-lokal einfach
zusammenh\"angend},
if each point $q \xeof \B$ a Umgebung $V_q$ besitzt
such that
each Schleife $\gamma$ at $q$ with $\gamma(S^1) \subset V_q$ is nullhomotop.

Alle for the Existenz einer universellen covering \textit{hinreichenden}
Eigenschaften of a topological space werden in dem
folgenden Begriff zusammengefa\"sst.

{\bf Definition 7} (hinreichend zusammenh\"angender space).

Ein topological space $\B$ is called \textit{hinreichend zusammenh\"angend},
if $\B$ wegzusammenh\"angend, lokal wegzusammenh\"angend
and semi-lokal einfach zusammenh\"angend ist.

\begin{example}
$S^1$ is hinreichend zusammenh\"angend.
\end{example}

Hinreichend zusammenh\"angende spaces besitzen the folgende Eigenschaft.

{\bf Satz 4}

Let $\B$ a hinreichend zusammenh\"angender topological space.
Then gibt es zu each subgroup $H \subset \pi_1(\B)$ eine
wegzusammenh\"angende covering $\pi:\E \to \B$
such that $\pi_1(\E) \cong H$.

Dieser Satz is the Grundlage for the Existenz einer
universellen \"Uberlagerung.

{\bf Satz 5} (Existenz einer universellen \"Uberlagerung).

Let $\B$ a hinreichend zusammenh\"angender topological space.
Then besitzt $\B$ a universal covering $\E$.

The universal covering nimmt a Sonderstellung unter allen
coverings of a space ein, and zwar in folgendem Sinne:
ist $\pi:\E \to \B$ a wegzusammenh\"angende \"Uberlagerung,
so \"uberlagert the universal covering von $\B$ auch $\E$.

Andererseits gilt:
if a universal covering von $\B$ existiert,
so is $\B$ semi-lokal einfach zusammenh\"angend.

\begin{example}
Let $\E \eq S^n$, $\B \eq \R\,P^n$ and $\pi(p) := \{p,-p\}$.
Im Falle $n \geq 2$ is $S^n$ einfach zusammenh\"angend and stellt damit
die universal covering von $\R\,P^n$ dar.
Im Falle $n=1$ is $S^1$ nicht einfach zusammenh\"angend and damit
\textit{not} a universal covering of $\R\,P^1$.
\end{example}

\section{Deck Transformations}
%------------------------------------------------------------------------------%
it Hilfe von \textit{Decktransformationen} kann man unter gewissen
Voraussetzungen the Fundamentalgruppen von topological spaces berechnen.

{\bf Definition 8} (Decktransformation).

Let $\pi:\E \to \B$ a \"Uberlagerung.
A \textit{Decktransformation} is a Hom\"oomorphismus $\phi:\E \to E$
such that
$\phi(\pi^{-1}(p)) \eq \phi(\pi^{-1}(p))$ for each $p \xeof \B$.

A Decktransformation f\"uhrt also Fasern wieder in Fasern \"uber.
The set $\cal{D}(\E)$ der Decktransformationen von $\E$
stellt with der Hintereinanderausf\"uhrung a group dar.

\begin{example}
$\R$ is a universal covering of $S^1$ and hence auch von
$\R\,P^1$.
The Decktransformationen are hier durch the Translationen
$x \mapsto x + 2k\pi$, $k \xeof \Z$ gegeben.
\end{example}

{\bf Satz 6}

Let $\pi:\E \to \B$ a covering and $\B$ wegzusammenh\"angend.
Then is the group der Decktransformationen isomorphic to a subgroup of $\pi_1(\B)$.

Folgende Eigenschaften des \"Uberlagerungsraumes bzw. des \"uberlagerten space
f\"uhren zur isomorphy der groups $\cal{D}(\E)$ and $\pi_1(\B)$.

{\bf Satz 7}

Let $\pi:\E \to \B$ a \"Uberlagerung.
Ist der topological space $E$ einfach zusammenh\"angend and
lokal wegzusammenh\"angend,
oder $\B$ hinreichend zusammenh\"angend and $E$ seine universal \"Uberlagerung,
then gilt $\cal{D}(\E) \cong \pi_1(\B)$.

\begin{examples}
\begin{itemize}
\item
Es gilt $\cal{D}(\R) \cong \Z$, also auch $\pi_1(S^1) \cong \Z$.
\item
Es gilt $\cal{D}(S^n) \cong \Z_2$ for $n \geq 2$.
The coverings $(S^n,\pi, \R\,P^n)$ erf\"ullen the Voraussetzungen
des obigen Satzes.
Es gilt daher $\pi_1(\R\,P^n) \cong \Z_2$ for $n \geq 2$ and
$\pi_1(\R\,P^1) \cong \pi_1(S^1) \cong \Z$.
\end{itemize}
\end{examples}

\section{Serre-Fibrations and Exact Homotopy-Sequences}
%------------------------------------------------------------------------------%
In point-set topology Das \textit{Hochhebeproblem} besteht darin,
unter welchen Voraussetzungen man a continuous mapping
$f:\topX \to \B$ zu einer continuous mapping $h:\topX \to \E$ hochheben kann
such that
for vorgegebenes $\pi:\E \to \B$ gilt: $\pi \circ h \eq f$.

{\bf Definition 9} (homotopy-Hochhebungseigenschaft).

A continuous mapping $\pi:\E \to \B$ hat the \textit{homotopy-Hochhebungs
eigenschaft} (HHE) in Bezug auf einen topological space $\topX$,
if es zu each continuous mapping $f:\topX \to \B$ a \textit{Hochhebung}
$h:X \to \E$ gibt, such that $f \eq \pi \circ h$.

\begin{picture}(200,90)
\put(180,75){$\E$}
\put(105,10){$\topX$}
\put(180,10){$\B$}
\put(120,12){\vector( 1, 0){55}}
\put(185,70){\vector( 0,-1){45}}
\put(120,15){\vector( 1, 1){57}}
\put(150,20){$f$}
\put(190,45){$\pi$}
\put(135,45){$h$}
\end{picture}

Folgende coverings are wie folgt ausgezeichnet:

{\bf Definition 10} (Serre-Faserung).

Let $I:=[0,1] \subset \R$.
A covering $\pi:\E \to \B$ is called a \textit{Serre-Faserung},
if sie the HHE for all topological spaces $I^n$, $n \xeof \N$ besitzt.

Gilt for a Serre-Faserung $\pi:\E \to \B$,
dass $\pi^{-1}(q) \simeq \F$ for einen topological space $\F$ and each point
$q \xeof \B$,
so sagt man, these besitze the \textit{typische Faser} $\F$.
Unter diesen Voraussetzungen gibt es a Instrument zur Berechnung der
h\"oheren homotopy groups,
n\"amlich the \textit{lange exakte homotopy-Sequenz einer Serre-Faserung}.

{\bf Satz 8}

Seien $\E,\B$ wegzusammenh\"angende topological spaces and
$(\E,\pi,\B)$ a Serre-Faserung with typischer Faser $\F$.
Let $p \xeof \B$ and $y \xeof \pi^{-1}(p)$ a beliebiger point.
Then is the Sequenz der homotopy groups
\[
\ldots \rightarrow \pi_n(\F,y) \rightarrow \pi_n(\E)
       \rightarrow \pi_n(\B)   \rightarrow
\ldots
\]
\[
\ldots \rightarrow \pi_1(\F,y) \rightarrow \pi_1(\E)
       \rightarrow \pi_1(\B)
\]
exakt.

Im Spezialfall, dass $\F$ diskret oder kontrahierbar ist,
ist $\pi_n(\F,y) \eq 0$.
Hieraus folgt, dass the Sequenzen
\[
0 \rightarrow \pi_n(\E) \rightarrow \pi_n(\B) \rightarrow 0
\]
exakt sind.
Es gilt also insbesondere $\pi_n(\E) \eq \pi_n(\B)$ for all $n \geq 2$.

\begin{example}
The mapping $\pi:\R \to S^1$ besitzt the diskrete Faser $\Z$,
also gilt $\pi_n(S^1) \eq 0$ for all $n \geq 2$.
\end{example}


\fancyhead[C] {\small{\textsc{IV. Chain Complexes}}}
%------------------------------------------------------------------------------%
\part{Chain Complexes}
%------------------------------------------------------------------------------%
Chain complexes treten vorwiegend bei der Untersuchung der globalen
Eigenschaften von topological Räumen auf.
Sie lassen sich rein algebraisch definieren and untersuchen.

\section{Short Exact Sequences of Modules}
%------------------------------------------------------------------------------%
Im folgenden werden Folgen von $R$-modules betrachtet,
die durch module homomorphisms miteinander verbunden sind. 
The homomorphisms sollen the Eigenschaft besitzen, 
dass the Verknüpfung zweier benachbarter homomorphisms stets the 
Nullabbildung ist.

{\bf Definition 1} (kurze Sequenz von $R$-modules)

Let $R$ a commutative ring with unit and $M_1,M,M_2$ jeweils
$R$-modules.
A \textit{kurze Sequenz} von $R$-modules is a Folge
\[
       \{0\}
       \longrightarrow M_1
       \stackrel{f}{\longrightarrow} M
       \stackrel{g}{\longrightarrow} M_2
       \longrightarrow               \{0\}
\]
with module homomorphisms $f:M_1 \to M$, $g:M \to M_2$ and
$g(f(x)) \eq 0$ for all $x \xeof M_1$.

Aufgrand der \textit{Komplexbedingung} $f \circ g \eq 0$ gilt stets the Inklusion 
$f(M_1) \subseteq \operatorname{Kern}(g)$.
Von besonderem Interesse are the cases, for which these sets are identical.

{\bf Definition 2} (kurze exakte Sequenz von $R$-modules)

A kurze Sequenz heißt \textit{exakt}, if $f(M) \eq \operatorname{Kern}(g)$.

\begin{example}
Seien $M_1,M_2$ jeweils $R$-modules and $f:M_1 \to M_2$ a homomorphism.
Then is the Sequenz
\[
       \{0\}
       \longrightarrow              \operatorname{Kern}(f)
       \longrightarrow               M_1
       \stackrel{f}{\longrightarrow} M_2
       \longrightarrow              \operatorname{Kokern}(f)
       \longrightarrow              \{0\}
\]
exakt.
Hieraus folgt, dass for a homomorphism $\phi:M_1 \to M_2$ gilt:
\begin{itemize}
\item
$\phi$ is genau then injektiv, if the Sequenz
$\{0\} \longrightarrow M_1 \stackrel{\phi}{\longrightarrow} M_2$ exakt ist.
\item
$\phi$ is genau then surjektiv, if the Sequenz
$M_1 \stackrel{\phi}{\longrightarrow} M_2 \longrightarrow \{0\}$ exakt ist.
\item
$\phi$ is genau then a Isomorphismus, if the Sequenz
$\{0\} \longrightarrow M_1 \stackrel{\phi}{\longrightarrow} M_2 
       \longrightarrow \{0\}$
exakt ist.
\end{itemize}
\end{example}

A kurze exakte Sequenz
\[
\{0\}     \longrightarrow M_1 \longrightarrow M 
\longrightarrow M_2 \longrightarrow    \{0\}
\]
von $R$-modules heißt \textit{spaltend},
if es Untermoduln $U_1,U_2$ von $M$ gibt
such that
$M \eq U_1 \oplus U_2$ and $f:M_1 \to U_1$ bzw. $g:U_2 \to M_2$ 
module-Isomorphismen sind.

\begin{picture}(100,60)
\put( 70,40){$M_1$}
\put( 95,50){$f$}
\put( 85,42){\vector(1,0){32}}
\put(120,40){$U_1$}
\put(120,25){$\bigoplus$}
\put(120,10){$U_2$}
\put(135,12){\vector(1,0){32}}
\put(170,10){$M_2$}
\put(145,20){$g$}
\put(110, 2){\line( 1, 0){30}}
\put(110, 2){\line( 0, 1){50}}
\put(140,52){\line(-1, 0){30}}
\put(140,52){\line( 0,-1){50}}
\end{picture}

The Begriff der kurzen Sequenz von modules läßt sich auf Folgen von modules
wie folgt erweitern.

{\bf Definition 3} (chain complex)

Let $R$ a commutative ring with unit.
\begin{itemize}
\item
Ein (absteigender) \textit{chain complex} $(K,d)$ is a Folge 
$\{M_i\}_{i \xeof \Z}$ von $R$-modules and module homomorphisms
$d_i: M_{i+1} \to M_{i}$,
die for all $i \xeof \Z$ der complex-condition
$d_{i} \circ d_{i+1} \eq 0$ genügen.
\item
Jede mapping $d_i$ heißt a \textit{Rand-homomorphism} of the chain complex.
\item
The sequence $\{d_i\}_{i \xeof \Z}$ der Rand-homomorphisms heißt das 
\textit{Differential} des chain complex.
\end{itemize}

Einen chain complex $(K,d)$ schreibt man gewöhnlich in der Form
\[
\ldots                    \longrightarrow  M_{i+1} 
       \stackrel{d_{i+1}}{\longrightarrow} M_i 
       \stackrel{d_{i}}  {\longrightarrow} M_{i-1} 
       \stackrel{d_{i-1}}{\longrightarrow} M_{i-2} 
       \longrightarrow 
\ldots
\]

\section{The Homology of a Chain Complex}
%------------------------------------------------------------------------------%
Wegen $d_{i} \circ d_{i+1} \eq 0$ gilt stets the Inklusion
$d_{i+1}(M_{i+1}) \subseteq \operatorname{Kern} (d_i)$.
Since the kernel of a module homomorphism s a $R$-module, too,
kann man for each $i \xeof \Z$ den Faktormodul
$\operatorname{Kern} (d_i) / d_{i+1}(M_{i+1})$ bilden.
Dies gibt Anlaß zu folgender Definition.

{\bf Definition 4} (Homologie-module, exakter chain complex)

Let $(K,d)$ a chain complex and $i \xeof \Z$.
\begin{itemize}
\item
The $R$-module
\[
H_i(K) := \frac {\operatorname{Kern} (d_i)}
                {d_{i+1}(M_{i+1})}
\]
heißt der $i$-te \textit{Homologie-module} des chain complex.
\item
$(K,d)$ heißt \textit{exakt} an der Stelle $i \xeof \N$,
if $H_i(K) \eq \{0\}$.
\item
Ein chain complex heißt \textit{exakt}, 
if er an each Stelle $i \xeof \Z$ exakt ist.
\end{itemize}

Hence the homology-modules of a chain complex messen the Abweichung
of the complex von seiner Exaktheit.

Ein Element $x \xeof M_{i}$ heißt a \textit{Zykel}, if $d_i(x) \eq 0$.
Gilt $x \eq d_{i-1}(y)$ for a $y \xeof M_{i-1}$, so heißt $x$ a \textit{Rand}.
Aufgrand der Komplex-Bedingung is each Zykel a Rand.
The Umkehrung gilt dagegen nur,
if der chain complex an der Stelle $i$ exakt ist.

Bildet man aus den homology-modules $H_i(K)$ the direkte sum,
so erhält man the \textit{totale homology} des chain complex:
\[
H(K) := \bigoplus_{i \xeof \Z} \, H_i(K).
\]
The totale homology $H(K)$ owns the structure of a $\Z$-graded $R$-module.

\section{Homomorphisms Between Chain Complexes}
%------------------------------------------------------------------------------%
Ist for zwei chain complexe $(K,d)$ and $(K',d')$ a Folge von 
module homomorphisms between den korrespondierenden modules
$M_i$ and $M'_i$ gegeben,
so kann man unter gewissen Zusatzvoraussetzungen a 
homomorphism between den chain complexes selbst erklären.

{\bf Definition 5} (Homomorphismus between chain complexes)

Seien $(K,d)$ and $(K',d')$ chain complexe.
A homomorphism $\phi: (K,d) \to (K',d')$ is a Folge 
$\{\phi_i\}_{i \xeof \Z}$ von module homomorphisms 
$\phi_i:M_i \to M'_i$
such that
\[
d'_i \circ \phi_i \eq \phi_{i-1} \circ d_i \qquad \mbox{for all } i \xeof \Z.
\]

For a homomorphism $\phi: (K,d) \to (K',d')$ erhält man also an each
Stelle $i \xeof \Z$ das folgende kommutative Diagramm:

\begin{picture}(100,100)
\put(140,15){$M'_{i}$}
\put(210,15){$M'_{i-1}$}
\put(140,75){$M_{i}$}
\put(210,75){$M_{i-1}$}
\put(162,77){\vector( 1, 0){45}}
\put(162,18){\vector( 1, 0){45}}
\put(148,70){\vector( 0,-1){45}}
\put(217,70){\vector( 0,-1){45}}
\put(133,48){$\phi_i$}
\put(223,48){$\phi_{i-1}$}
\put(175,80){$d_i$}
\put(175,25){$d'_{i}$}
\end{picture}

The chain complexes together with their homomorphisms constitute a category.

\section{Homotopy of Chain Complexes}
%------------------------------------------------------------------------------%
Ist $\phi$ a homomorphism between den chain complexes $(K,d)$ and 
$(K',d')$,
so wird hierdurch a Folge $\{\phi^\*_i\}_{i \xeof \Z}$
von module homomorphism between den homology-modules der 
beiden chain complexe induziert: 
\[
\phi^\*_i : H_i(K) \to H_i(K'), \quad \phi_i([x]) := [\phi_i(x)].
\]

Seien $\phi$, $\psi$ zwei homomorphisms between den chain complexes
$(K,d)$ and $(K',d')$.

{\bf Definition 6} (zueinander homotope chain complexe)

Seien $(K,d)$ and $(K',d')$ chain complexe.
Zwei homomorphisms $\phi, \psi: (K,d) \to (K',d')$ are called
\textit{zueinander homotop},
if for each $i \xeof \Z$ a module-homomorphism
$\alpha_i:M_i \to M'_{i+1}$ existiert
such that
\[
\phi_i - \psi_i \eq d'_{i+1} \circ \alpha_i + \alpha_{i-1} \circ d_i.
\]

Folgendes Diagramm veranschaulicht these Identitäten.

\begin{picture}(100,100)
\put( 70,15){$M'_{i+1}$}
\put(145,15){$M'_{i}$}
\put(210,15){$M'_{i-1}$}
\put( 70,75){$M_{i+1}$}
\put(145,75){$M_{i}$}
\put(210,75){$M_{i-1}$}
\put( 92,77){\vector( 1, 0){45}}
\put( 92,18){\vector( 1, 0){45}}
\put(162,77){\vector( 1, 0){45}}
\put(162,18){\vector( 1, 0){45}}
\put( 78,70){\vector( 0,-1){45}}
\put(148,70){\vector( 0,-1){45}}
\put(217,70){\vector( 0,-1){45}}
\put( 55,45){$\phi_{i+1}$}
\put( 55,55){$\psi_{i+1}$}
\put(150,45){$\phi_i$}
\put(150,55){$\psi_i$}
\put(223,45){$\phi_{i-1}$}
\put(223,55){$\psi_{i-1}$}
\put(105,82){$d_{i+1}$}
\put(175,82){$d_i$}
\put(105,25){$d'_{i+1}$}
\put(180,25){$d'_{i}$}
\put(140,70){\vector( -1,-1){47}}
\put(210,70){\vector( -1,-1){47}}
\put(168,50){$\alpha_{i-1}$}
\put(108,50){$\alpha_{i}$}
\end{picture}

homotopy is a Äquivalenzrelation auf der set der homomorphisms 
between zwei chain complexes.
Zwei homomorphisms aus derselben Äquivalenzklasse are called zueinander
\textit{kettenhomotop}.

{\bf Definition 7} (homotopie-äquivalente chain complexe)

Zwei chain complexe $(K,d)$ and $(K',d')$ are called \textit{homotopie-äquivalent},
wenn es homotope homomorphism 
$\phi:(K,d)   \to (K',d')$ and 
$\psi:(K',d') \to (K,d)$
gibt, such that $\phi \circ \psi \eq id_{K'}$ bzw. $\psi \circ \phi \eq id_K$.
The homomorphisms $\phi$ and $\psi$ are called in diesem Falle 
\textit{zueinander homotopie-invers}.

The Bedeutung der homotopy-Äquivalenz wird durch folgenden Satz ersichtlich.

{\bf Satz 1}

Seien $(K,d)$ and $(K',d')$ homotopie-äquivalente chain complexe.
Then are the induzierten module-homomorphisms between den korrespondierenden
Homologie-modules zueinander isomorph:
\[
f_{\*} \eq g_{\*} : H_i(K) \to H_i(K').
\]

Für homotopie-inverse homomorphisms schreibt man kurz $\phi \simeq \psi$.
The homotopy-Äquivalenz is a Äquivalenzrelation
in der category der chain complexe.

The Begriff der kurzen exakten Sequenz von modules kann auf chain complexe
wie folgt erweitert werden.

{\bf Definition 8} (kurze exakte Sequenz von chain complexes)

Ein \textit{kurze exakte Sequenz} von chain complexes is a Folge
\[
\{0\} 
\longrightarrow K' 
\stackrel{f}{\longrightarrow} K 
\stackrel{g}{\longrightarrow} K''
\longrightarrow 
\{0\}
\]
von Komplexen and homomorphisms
such that for each $i \xeof \Z$ the Folge der modules
\[
\{0\} \longrightarrow M'_i 
      \stackrel{f_i}{\longrightarrow} M_i 
      \stackrel{g_i}{\longrightarrow} M''_i 
      \longrightarrow
\{0\}
\]
exakt ist.

Eines der wichtigsten Resultate in der homologischen Algebra ist,
dass a kurze exakte Sequenz von chain complexes
eine lange exakte Sequenz von homology-modules erzeugt.

{\bf Satz 2} (Satz von der langen exakten homology-Sequenz).

Sei
$
\{0\} 
\longrightarrow K' 
\stackrel{f}{\longrightarrow} K 
\stackrel{g}{\longrightarrow} K''
\longrightarrow 
\{0\}
$
eine kurze exakte Sequenz von chain complexes.
Then is for all $n \xeof \N$ the mapping
$\delta_n: H_n(K'') \to H_{n-1}(K')$ durch the mappingsvorschrift
$[z] \xeof H_n(K'') \mapsto [ f^{-1}_{n-1} \circ d_n \circ g^{-1}_n(z) ]$
wohldefiniert and folgende lange Sequenz exakt:
\[
\ldots 
       \stackrel{\delta_{n+1}}{\longrightarrow} H_n(K')
       \stackrel{ f_n}        {\longrightarrow} H_n(K)
       \stackrel{ g_n}        {\longrightarrow} H_n(K'')
       \stackrel{\delta_n}    {\longrightarrow} H_{n-1}(K')
\ldots
\]

The mappings $\delta_n$ are called the \textit{Randhomomorphismen} der
kurzen exakten Sequenz von chain complexes.

The direkte sum von chain complexes ergibt sich in naheliegender Weise
aus der direkten sum von modules.

{\bf Definition 9} (direkte sum von chain complexes)

Seien $(K',d')$ and $(K'',d'')$ zwei chain complexe.
The \textit{direkte sum} von $(K',d')$ and $(K'',d'')$ is der chain complex
$(K,d)$, der durch $M_i := M'_i \oplus M''_i$ and 
$d_i(x,y) := (d'(x),d''(y))$
for $x \xeof M'_i$ and $y \xeof M''_i$ definiert ist.

Let $(K' \oplus K'', d' \oplus d'')$ the direkte sum der chain complexe
$(K',d')$ and $(K'',d'')$.
The homology-modules dieser direkten sum lassen sich
aus den homology-modules der Summanden berechnen:
\[
H_i(K' \oplus K'') \cong H_i(K') \oplus H_i(K'') \quad \fa i \xeof \Z.
\]

\section{Sub-Complexes of Chain Complexes}
%------------------------------------------------------------------------------%
an kann in naheliegender Weise von Unterkomplexen of a chain complex
sprechen,
wenn for each module $M_i$ of the complex a Untermodul $U_i$ vorgegeben ist
such that
the Einschränkungen der Randhomomorphismen $d_i$ auf die
Untermoduln $U_i$ also Randhomomorphismen sind
and the Komplex-Bedingung $d_i \circ d_{i+1} \eq 0$ erhalten bleibt.

{\bf Definition 10} (Unterkomplex)

Let $(K,d)$ a chain complex and $\{U_i\}_{i \xeof \Z}$ for each 
$i \xeof \Z$ a Untermodul von $M_i$, such that $d_i(U_i) \subseteq U_{i-1}$.
Then heißt $(U,d)$ a \textit{Unterkomplex} von $(K,d)$.

Let $(U,d)$ a Unterkomplex of a chain complex $(K,d)$.
Bildet man for each $i \xeof \Z$ the Faktor-modules $M_i/U_i$,
so erhält man wieder a Folge von (Faktor)-modules. 
Für these definiert man module-homomorphisms 
\[
d''[x]: M_i/U_i \to M_{i-1}/U_{i-1}, \quad d''[x] := [d(x)]
\]
such that
with $M''_i := K_i/U_i$ the Folge $(K'',d'')$ wieder einen
chain complex bildet.

{\bf Definition 11} (Faktorkomplex)

The oben angeführte chain complex $(K'',d'')$ heißt der Faktorkomplex bzw. 
Quotienten-Komplex des chain complex $(K,d)$ and des Unterkomplexes 
$(U,d)$.

Let nun $(K,d)$ a chain complex, $(U,d)$ a Unterkomplex and $(K'',d'')$
der hieraus gebildete Faktorkomplex. 
it den Einbettungsabbildungen $\iota_i:U_i \hookrightarrow M_i$
and den Projektionsabbildungen $\pi_i: M_i \to M''_i$
stellen the drei Komplexe a kurze exakte Sequenz von chain complexes dar.
Es existiert daher zu $(U,d)$, $(K,d)$ and $(K'',d'')$ a lange exakte 
Sequenz von homology-modules and module-homomorphisms.

Bildet man for zwei Unterkomplexe modulweise the intersection of the modules
$M'_i, M''_i$, so erhält man auf these Weise den chain complex 
$K' \cap K''$.

{\bf Satz 3}

Let $(K,d)$ be a chain complex and $(K',d'), (K'',d'')$ subcomplexes.
Seien ferner the mappings $i$ and $p$ defined according to
$i(x) \mapsto (x,-x)$ for $x \xeof K'$ 
and $p(x,y) \mapsto x+y$ for $(x,y) \xeof K' \oplus K''$.
Then gilt: the kurze Sequenz von chain complexes
\[
0              \longrightarrow  K' \cap   K'' 
  \stackrel{i}{\longrightarrow} K' \oplus K''
  \stackrel{p}{\longrightarrow} K
               \longrightarrow
0
\]
ist exakt. 

\section{The Tensor Product of Chain Complexes}
%------------------------------------------------------------------------------%
Es läßt sich sowohl a Tensorprodukt of a chain complex with einem
$R$-module als auch with einem weiteren chain complex definieren.

{\bf Definition 12} (Tensorprodukt von chain complexes)

Let $R$ a commutative ring with unit,
$(K,d)$ a chain complex and $M$ a $R$-module.
Then is das Tensorprodukt von $(K,d)$ with $M$ der chain complex 
$(\tensor{K}{M}, c)$ with den modules $\tensor{K_i}{M}$ and den
homomorphisms $c_i: \tensor{K_i}{M} \to \tensor{K_{i-1}}{M}$,
$c_i := d_i \otimes id_M$.

Let $H_i(K,M) := H_i(\tensor{K}{M},c)$.
The modules $H_i(K,M)$ are called the homology-modules of the chain complex 
$(K,d)$ with coefficients in $M$.

The tensor product of two chain complexe is defined as follows:

{\bf Definition 12} (Tensorprodukt von chain complexes)

Let $R$ a commutative ring with unit and $(K,d)$, $(L,d')$ 
chain complexe.
Then is das Tensorprodukt $\tensor{(K,d)}{(K',d')}$ der
chain complex $(C,e)$ mit
\[
C_n := \bigoplus_{p+q=n} \tensor{K_p}{{K'}_q}
\]
und
\[
e_n(\tensor{x_p}{y_q}) := \tensor{d_p(x_p)}{y_q} + 
             (-1)^p \cdot \tensor{x_p}{d'_q(y_q)},
\qquad x_p \xeof K_p, \, y_q \xeof {K'}_q.
\]

\section{Cohomology of Chain Complexes}
%------------------------------------------------------------------------------%
Let $(K,d)$ a absteigender chain complex.
Hieraus erhält man in a natural way einen \textit{aufsteigenden}
chain complex,
indem man the modules and module-homomorphisms umnummeriert:
$K^n := K_{-n}$ and $d^n := d_{-n}$ for $n \xeof \Z$.

\subparagraph{}
Durch Anwenden des Hom-functors $Hom(-,M)$ auf einen absteigenden chain complex
$(K,d)$ erhält man also einen aufsteigenden chain complex,
and zwar in folgender Weise:
zum homomorphism $d_n:K_n \to K_{n-1}$ sei
\[
d^{\*}_n:Hom_R(K_{n-1},M) \to Hom_R(K_n,M)
\]
der \textit{duale} homomorphism sowie $d^n := (-1)^n \, d^{\*}_{n+1}$.
In diesem Falle heißt $(Hom(K,M),d)$ a \textit{cochain complex}.

{\bf Definition 13} (cochain module, Korand, Kozykel)

Let $(K,d)$ a chain complex, $M$ a $R$-module and $(Hom(K,M),d)$ der
zugehörige cochain complex.
Then heißt $Hom_R(K_n,M)$ der $n$-te cochain module von $(K,d)$
with values in $M$.
Ein Element aus $\operatorname{Kern}(d^n)$ heißt \textit{$n$-Kozykel},
and a Element von $d^{n-1}(K^{n-1})$ heißt \textit{$n$-Korand} von
$K$ with coefficients in $M$.

\subparagraph{}
Analog zu den homology-modules definiert man:

{\bf Definition 14} (Kohomologie-module)

Let $R$ a commutative ring with unit and $(K,d)$ der aufsteigender
chain complex
\[
\ldots
\stackrel{d^{n-2}}{\longrightarrow}
K^{n-1}
\stackrel{d^{n-1}}{\longrightarrow}
K^{n}
\stackrel{d^{n}}{\longrightarrow}
K^{n+1}
\stackrel{d^{n+1}}{\longrightarrow}
\ldots
\]
Then heißt
$H^n(K) := \operatorname{Kern}(d^n) / d^{n-1}(K^{n-1})$
der $n$-te \textit{Kohomologie-module} von $(K,d)$.

\fancyhead[C] {\small{\textsc{V. Simplicial Complexes}}}
%------------------------------------------------------------------------------%
\part{Simplicial Complexes}
%------------------------------------------------------------------------------%
\textit{Simpliziale Komplexe} bzw. \textit{Polyeder} are subsets des $\R^n$,
die aus sogenannten \textit{Simplices} zusammengesetzt sind.
Viele topological spaces k\"onnen durch simplicial complexes
beliebig approximiert werden.

\section{Simplicial Chains}
%------------------------------------------------------------------------------%
A Anzahl von points $p_0,\ldots,p_k \xeof \R^n$ with $k \leq n$ are called
\textit{in allgemeiner Lage},
if for $i=1,\ldots,k$ the Vektoren $v_i := p_i-p_0$ linear unabh\"angig sind,
and in diesem Fall is the \textit{konvexe H\"ulle} $[p_0,\ldots,p_k]$ erkl\"art
als the set
\[
[v_1,\ldots,v_k]
\, := \,
\set{x \xeof \R^n}
    {x \eq p_0 \, + \, \sum_{i=1}^k \lambda_i \, v_i,
    \quad 
    \lambda_i \geq 0,
    \quad
    \lambda_1 + \ldots + \lambda_k \leq 1}.
\]
For einen point $x \xeof [p_0,\ldots,p_k]$ are called the coefficients
$\lambda_i$ the \textit{baryzentrischen Koordinaten} von $x$.
Dawith kann der Begriff \textit{Simplex} wie folgt erkl\"art werden:

{\bf Definition 1} ($k$-dimensionales Simplex).

Ein \textit{$k$-dimensionales Simplex} bzw. \textit{$k$-Simplex} $\simS \subset \R^n$
ist the konvexe H\"ulle von $k+1$ points $p_i$ in allgemeiner Lage.
Ein \textit{Teilsimplex} von $\simS$ is the konvexe H\"ulle einer subset von
$\{ p_0,\ldots,p_k\}$.

\begin{examples}
\begin{itemize}
\item
Jeder point $p \xeof \R^n$ is a $0$-Simplex and wird a \textit{Ecke} genannt.
\item
The $1$-Simplices are Segmente with Endpunkten $p,q \xeof \R^n$ \textit{(Kanten)}.
\item
The $2$-Simplices are nichtentartete Dreiecke,
bestehend aus jeweils drei Segmenten,
die wiederum Teil-Simplices hiervon sind.
\end{itemize}
\end{examples}

The \textit{Rand} bzw. das \textit{Innere} of a Simplices is wie folgt erkl\"art:

{\bf Definition 2} (Rand of a Simplices, offener Simplex).

Let $\simS$ a $k$-Simplex.
The union $\partial \, \simS$ der Teilsimplices von $\simS$ is called der
\textit{Rand} von $\simS$.
The set $S^0 := \simS \setminus \partial \, \simS$
is called der \textit{offene Simplex} von $\simS$.

\begin{example}
Ecken besitzen als Rand the leere set,
each Kante hat seine beiden Endpunkte als Rand,
and der Rand of a $2$-Simplices besteht aus seinen drei Kanten.
\end{example}

Let $\simS$ a $k$-dimensionales Simplex,
der von den $k+1$ points $p_0,\ldots,p_k$ in allgemeiner Lage erzeugt wird.
Then is called der point
\[
\overline{\simS} \, := \, \frac{1}{k+1} \, \sum_{i=0}^k\,  p_i \quad \xeof \R^n
\]
der \textit{Schwerpunkt} von $\simS$.
The nichtnegative reelle Zahl
$\delta_{\simS} := \operatorname{sup} \set{\norm{x-y}}{x,y \xeof \simS}$
is called der \textit{Durchmesser} of a Simplices.

\section{Polyeders}
%------------------------------------------------------------------------------%
Baut man Simplices so zusammen,
dass points, edges, Dreiecke usw. zur Deckung gebracht werden,
so erh\"alt man einen sogenannten \textit{simplizialen Komplex}.

{\bf Definition 3} (simplizialer Komplex, Polyeder).

A set $\simK$ von Simplices is called a \textit{simplizialer Komplex}
bzw. a \textit{Polyeder}, if gilt:
\begin{itemize}
\item
it each Simplex $\simS \xeof \simK$ enth\"alt $\simK$ auch
s\"amtliche Teilsimplices von $\simS$.
\item
The intersection of je zwei Simplices $\simS_1$, $\simS_2 \xeof \simK$
ist entweder leer oder a gemeinsamer Teilsimplex.
\item
Enth\"alt $\simK$ unendlich viele Simplices,
so is $\simK$ \textit{lokal endlich},
that is,
each point $x \xeof \R^n$ hat a Umgebung,
die nur with endlich vielen Simplices aus $K$ nichtleeren Durchschnitt besitzt.
\end{itemize}

Let $\simK$ a simplizialer Komplex.
Then is the Dimension von $\simK$
das Maximum der Dimensionen aller Simplices von $\simK$.

A subset $\simK'$ of a simplizialen Komplexes $\simK$ is called ein
\textit{Unterkomplex} von $\simK$,
if $\simK'$ selbst einen simplizialen Komplex darstellt.
The Dimension of a Unterkomplexes is offensichtlich h\"ochstens so gro\"ss
wie the Dimension of the complex selbst.

Let $ \abs{\simK} $ the mengentheoretische union aller
Simplices of a simplizialen Komplexes $\simK$:
\[
\abs{\simK} \, := \, \bigcup_{\simS_i \xeof \simK} \simS_i.
\]
Then is called $\abs{ \simK } \subset \R^n$ der dem Komplex zugrundeliegende
topological space.
The Topologie von $ \abs{\simK} $ is dabei durch the Spur-Topologie
des $\R^n$ in $ \abs{\simK} $ vorgegeben.
Zu each $p \xeof \abs{\simK} $ gibt es genau a Simplex $\simS_p \xeof \simK$
with maximaler Dimension and $p \xeof \simS_p$.
Dieses is called das \textit{Tr\"ager-Simplex} von $p$.

Ist andererseits $\topX$ a beliebiger topological space
and $\simK$ a (endlicher) simplizialer Komplex,
so is called $\topX$ durch $\simK$ \textit{(endlich) triangulierbar},
if $\topX$ zu $ \abs{\simK} $ hom\"oomorph ist.

\begin{example}
The Kreisscheibe $D^1$ is triangulierbar,
da hom\"oomorph zu einem Dreieck $\Delta \subset \R^2$.
\end{example}

\section{Simplicial Mappings}
%------------------------------------------------------------------------------%
\textit{Simpliziale mappings} are continuous mappings
between simplizialen Komplexen,
die Simplices wiederum in Simplices \"uberf\"uhren.

{\bf Definition 4} (simpliziale mapping).

Seien $\simK$ and $\simL$ simplicial complexes.
A \textit{simpliziale mapping} $f:\simK \to \simL$ ist
eine continuous mapping with folgenden Eigenschaften:
\begin{itemize}
\item
$f$ bildet the Ecken von $\simK$ in Ecken von $\simL$ ab.
\item
$f(\simS) \xeof \simL$ for each Simplex $\simS \xeof \simK$.
\end{itemize}

A simpliziale mapping is eindeutig durch the Bilder der Ecken
bzw. der $0$-Simplices bestimmt.
Sind $\simK$ and $\simL$ zwei simplicial complexes
with jeweils $m$ bzw. $n$ Ecken,
so gibt es h\"ochstens $n^m$ simpliziale mappings $f:\simK \to \simL$
(meistens are es jedoch weniger).

For a simpliziale mapping $f$ gilt stets:
$\dim (\simS) \geq \dim (f(\simS))$ for each Simplex $\simS \xeof \simK$.

A simpliziale mapping $f$ between den Komplexen $\simK$ and $\simL$
induziert a continuous mapping $ \abs{f} $ between den topological spaces
$ \abs{\simK} $ and $ \abs{\simL} $.
The mapping $ \abs{f} $ is for einen point $p \xeof \abs{\simK} $ with den
baryzentrischen Koordinaten $\lambda_i$ bez\"uglich des Tr\"ager-Simplices
$(x_0,\ldots,x_n)$ erkl\"art according to
\[
\abs{f}(p) \, := \, \sum_{i=0}^n \lambda_i \, f(x_i).
\]
Ist $f:\simK \to \simL$ a simpliziale mapping and \"uberdies bijektiv,
so is auch $f^{-1}$ simplizial.
Insbesondere is then $ \abs{f} $ a Hom\"oomorphismus between
den spaces $ \abs{\simK} $ and $ \abs{\simL} $.

The simplizialen Komplexe bilden with zusammen with den simplizialen
mappings a category.

\section{The Simplicial Approximation Theorem}
%------------------------------------------------------------------------------%
Das Ziel der \textit{simplizialen Approximation} is es,
eine beliebige continuous mapping between Polyedern durch a simpliziale
mapping anzun\"ahern.
Hierzu m\"ussen in der Regel the Polyeder verfeinert werden.

{\bf Definition 5} (baryzentrische Unterteilung).

Let $\simK$ a simplizialer Komplex.
A \textit{baryzentrische Unterteilung} von $\simK$ ist
ein weiterer simplizialer Komplex $\simK'$,
dessen Ecken the Schwerpunkte von Simplices in $\simK$ sind.

Da each Simplex von $\simK$ hierbei in kleinere Simplices
von $\simK'$ zerlegt wird,
ist the baryzentrische Unterteilung von $\simK$ tats\"achlich a Verfeinerung.
Hierbei gilt for the Durchmesser von Simplices von $\simK'$:
\[
d(S') \leq \frac{\dim(\simK)}{\dim(\simK)+1}
\cdot
\operatorname{max} \set{d(S)}{S \xeof \simK}.
\]
The zugrundeliegende topological space bleibt bei einer Verfeinerung erhalten,
that is,
es gilt $ \abs{\simK} \eq \abs{\simK'} $.

Das Verfahren der baryzentrischen Unterteilung kann iterativ angewandt werden.
Let $K^0 := K$ and $K^n := (K^{n-1})'$ the $n$-te baryzentrische Unterteilung
von $\simK$.
Zu $\epsilon > 0$ gibt es $n \xeof \N$
such that
each simplex of $K^n$ owns a diameter,
which is smaller than $\epsilon$.

{\bf Definition 6} (simpliziale Appoximation).

Seien $\simK$ and $\simL$ simplicial complexes and
$f: \abs{\simK} \to \abs{\simL} $ a continuous mapping.
A simpliziale mapping $\phi:\simK \to \simL$ is called eine
\textit{simpliziale Approximation} von $f$,
wenn for all $p \xeof \abs{\simK} $ gilt:
$\abs{\phi}(p) \xeof \overline{S_{f(p)}}$.

The folgende Satz von der \textit{Existenz einer simplizialen Approximation}
hat zum Inhalt,
dass man each continuous mapping between Polyedern
durch a simpliziale mapping approximieren kann,
wenn man den Urbildraum durch baryzentrische Unterteilungen
hinreichend verfeinert.

{\bf Satz 1} (Existenz einer simplizialen Approximation).

Seien $\simK$ and $\simL$ simplicial complexes
and $f: \abs{\simK} \to \abs{\simL} $ a continuous mapping.
Then gibt es a $n_0 \xeof \N$ and for each $n \geq n_0$ a simpliziale
Approximation $\phi_n:K^n \to \simL$ von $f$.

\section{Ordered Simplicial Complexes}
%------------------------------------------------------------------------------%
Um a homology-Theorie for simplicial complexes bzw. ihre assoziierten
topological spaces formulieren zu k\"onnen,
ben\"otigt man a Totalordnung auf der set der $0$-Simplices of the complex.

{\bf Definition 7} (geordneter simplizialer Komplex).

Let $\simK$ a finite simplicial complex.
$\simK$ is called \textit{geordnet},
if auf der set aller $0$-Simplices a Totalordnung $<$ erkl\"art ist.

The points of a geordneten Komplexes are also durchnummeriert.
Let $(\simK,<)$ the Bezeichnung for einen \textit{geordneten} simplizialen Komplex.
Then l\"a\"sst sich a $k$-Simplex $\simS \xeof \simK$ eindeutig in der Form
$\simS \eq [p_0,\ldots,p_k]$ darstellen.

{\bf Definition 8} ($i$-te Seite of a geordneten simplizialen Komplexes).

Let $(\simK,<)$ a geordneter simplizialer Komplex.
For a $k$-Simplex $\simS \eq [p_0,\ldots,p_k] \xeof \simK$ and $0 \leq i \leq k$
is called das $(k-1)$-Simplex $\partial_i \simS$ mit
$\partial_i \simS := \, [p_0,\ldots,p_{i-1},p_{i+1},\ldots,p_k]$
die $i$-te Seite von $\simS$.

The Seite $\partial_i \simS$ is also dasjenige Teilsimplex von $\simS$,
das dem point $p_i$ gegen\"uber liegt.

Let $K_n$ diejenige subset des simplizialen Komplexes $\simK$,
die genau the $n$-Simplizes enth\"alt.
Then f\"uhrt the Zuordnung $\simS \mapsto \partial_i \simS$
Simplices aus $K_n$ in Simplices von $K_{n-1}$ \"uber.

\section{The Euler-Charakteristics of a Simplicial Complex}
%------------------------------------------------------------------------------%
Let im folgenden $\riR$ a commutative ring with unit,
$(\simK,<)$ a geordneter simplizialer Komplex and $n \xeof \N$.
The set $C_n(\simK,R)$ sei der von den $n$-Simplices
frei erzeugte $\riR$-module,
that is,
each element $c \xeof C_n(\simK,R)$ is a formal sum of the form
\[
c \, \eq \, \sum_{i=1}^n r_i K_i, \quad r_i \xeof R.
\]

{\bf Definition 9} (simpliziale $n$-chain).

The elements von $C_n(\simK,R)$ are called \textit{simpliziale $n$-chains}
in $\simK$ (with coefficients in $\riR$).

Let for $n > 0$ the mapping $d_n: C_n(\simK,R) \to C_{n-1}(\simK,R)$
definiert according to
\[
d_n(S) := \sum_{i=0}^n (-1)^i \, \partial_i \, S, \qquad S \xeof C_n(\simK,R).
\]
Then stellt the mapping $d_n$ for each $n \xeof \N$
einen module-homomorphism dar.

\begin{example}
Sind $k_0, k_1, k_2$ the Kanten of a Dreiecks $\Delta$,
so gilt $d_2(\Delta) \eq k_0 - k_1 + k_2$ and $d_1 \circ d_2(\Delta) \eq 0$.
\end{example}

Es is a wichtige Eigenschaft von chains,
dass the zweimalige Anwendung der Randbildung einer chain
die $0$-chain liefert:

{\bf Satz 2}

Es gilt $d_{n-1} \circ d_{n} \eq 0$ for $n \geq 2$.

The sequence der $\riR$-modules $C_n(\simK,R)$ stellt also
with den homomorphisms $d_n$ a chain complex $(C(\simK,R),d)$ dar.
Es is naheliegend, the homology-modules dieses chain complex untersuchen:

{\bf Definition 10} (simplizialer homology-module).

Let $(\simK,<)$ a geordneter simplizialer Komplex.
Then is called $H_n(\simK;\riR)$
der $n$-te \textit{simpliziale Homologie-Modul} von $(\simK,<)$.

it $H_{\*}(\simK,R)$ wird gew\"ohnlich the \textit{totale simpliziale homology}
des geordneten simplizialen Komplexes $(\simK,<)$ bezeichnet.

Zwischen den homology-modules and der \textit{Euler-Charaktistik} of the complex,
die historisch gesehen the \"alteste topological Invariante darstellt,
gibt es einen engen Zusammenhang.

{\bf Definition 11} (Euler-Charakteristik).

Let $\simK$ a endlicher simplizialer Komplex and $a_n(\simK)$ the Anzahl
der $n$-Simplices von $\simK$.
Then is called the Zahl
\[
\chi(\simK) := \sum_{n=0}^{\infty} a_n(\simK) \, \xeof \Z
\]
die \textit{Euler-Charakteristik} von $\simK$.

The \textit{Eulersche Polyedersatz} liefert for den Fall,
dass $\riR$ a K\"orper ist,
folgende Beziehung between der Euler-Charakteristik of a Komplexes and
seinen homology-modules.

{\bf Satz 3} (Eulerscher Polyeder-Satz).

Let $(\simK,<)$ a finite and ordered simplicial complex and $\riR$ a field.
Then we have
\[
\chi(\simK) \eq \sum_{n=0}^{\infty} \dim_{\riR} \, H_n(\simK;\riR).
\]

The values $\beta_n := \dim_{\riR} \, H_n(\simK;\riR)$ are called
die \textit{Betti numbers} of the complex.

\begin{example}
The Betti numbers of a W\"urfels $W$ are $\beta_0 \eq 8$, $\beta_1 \eq 12$,
$\beta_2 \eq 6$ and $\beta_3 \eq 1$,
which implies $\chi(W) \eq 8-12+6-1 \eq 1$.
\end{example}

\section{CW-Complexes}
%------------------------------------------------------------------------------%
CW-Komplexe are Verallgemeinerungen von simplizialen Komplexen.
Viele topological spaces wie z.B. Sph\"aren oder Tori k\"onnen with CW-Komplexen
wesentlich einfacher approximiert werden als with Polyedern.

\textit{cells} and Zerlegungen von topological spaces in cells bilden die
Grundlagen der CW-Komplexe.

{\bf Definition 12} ($n$-Zelle).

Let $\topX$ a topological space.
\begin{itemize}
\item
The space $\topX$ is called a \textit{$n$-Zelle},
if er zum space $\R^n$ hom\"oomorph ist.
\item
Jeden point $p$ einer $n$-Zelle is called a $0$-Zelle.
\item
A \textit{cellszerlegung} $(\cal{E},\topX)$ von $\topX$ is a setnsystem
$\cal{E}$ von $n$-cells,
die paarweise disjunkt are and den space $\topX$ \"uberdecken.
\end{itemize}

\begin{examples}
\begin{itemize}
\item
The offene Vollkugel $D^n$ sowie the punktierte Sph\"are
$S^n \setminus \{p\} \subset \R^{n+1}$ are $n$-cells.
\item
The open simplices of a simplizialen Komplexes $\simK$ are $n$-cells,
and all open simplices of the complex bilden a cellszerlegung
des topological space $ \abs{\simK} $.
\end{itemize}
\end{examples}

The cellszerlegungen of a space sollten gewissen Einschr\"ankungen
unterworfen sein,
die zum Begriff des \textit{CW-Komplexes} f\"uhren:

{\bf Definition 13} (CW-Komplex).

Let $\topX$ a Hausdorffscher topological space and $\cal{E}$
eine cellszerlegung von $\topX$.
Then is called $(\cal{E},\topX)$ a \textit{CW-Komplex}, if gilt:
\begin{itemize}
\item
Zu each $n$-Zelle $\epsilon \xeof \cal{E}$ gibt es a continuous mapping
$\Phi:D^n \to \topX$,
die das Innere von $D^n$ hom\"oomorph auf the Zelle $\epsilon$
and $S^{n-1}$ in the union der h\"ochstens $(n-1)$-dimensionalen cells
abbildet \textit{(charakteristische mapping)}.
\item
The abgeschlossene H\"ulle $\overline{\epsilon}$ each Zelle
$\epsilon \xeof \cal{E}$ hat nur with endlich vielen anderen cells
non-empty intersection \textit{(H\"ullen-Endlichkeit)}.
\item
A subset $\setM \subset \topX$ is genau then abgeschlossen,
wenn each set $\setM \cap \overline{\epsilon}$ abgeschlossen ist
\textit{(schwache Topologie)}.
\end{itemize}

Ist a cellszerlegung endlich,
so is the Forderung nach der H\"ullen-Endlichkeit
and der schwache Topologie bereits erf\"ullt.

\begin{examples}
\begin{itemize}
\item
Jeder simpliziale Komplex is a CW-Komplex.
\item
The spaces $S^n$ are CW-Komplexe.
Insbesondere besteht the Sph\"are $S^2$ aus einer $0$-Zelle (Nordpol) and
einer $2$-Zelle.
\item
The Torus is a CW-Komplex, bestehend aus einer $0$-Zelle, zwei $1$-cells
and einer $2$-Zelle.
\item
The $3$-dimensionale W\"urfel besteht aus acht $0$-cells, zw\"olf $1$-cells,
sechs $2$-cells and einer $3$-Zelle.
\item
The projektive space $\R P^n$ besitzt a CW-Zerlegung
\[
\R P^0 \subset \R P^1 \subset \ldots \subset \R P^n
\]
with einer $i$-Zelle for each $i=0,\ldots,n$.
\item
The projektive space $\C P^n$ besitzt a CW-Zerlegung
\[
\C P^0 \subset \C P^1 \subset \ldots \subset \C P^n
\]
with einer $2i$-Zelle for each $i=0,\ldots,n$.
\end{itemize}
\end{examples}

CW-Komplexe are bez\"uglich der performing the intersection and the union abgeschlossen.

The Aufbau von CW-Komplexen erfolgt gew\"ohnlich with Hilfe von \textit{$n$-Ger\"usten},
welche folgenderma\"ssen definiert sind:

{\bf Definition 14} ($n$-Ger\"ust).

Let $(\cal{E},\topX)$ a CW-Komplex.
For $n \xeof \N$ is called the set $\cal{E}^n$ aller $n$-cells $\epsilon
\in \cal{E}$ das $n$-Ger\"ust bzw. $n$-Skelett von $(\cal{E},\topX)$.

Let $\epsilon$ a $n$-Zelle.
The Restriktion der charakteristischen mapping $\Phi:S^{n-1} \to \cal{E}^n$
is called the \textit{anheftende mapping} der Zelle an das $n-1$-Ger\"ust.
The $n$-cells von $\cal{E}$ are also durch $\Phi$ an das $n-1$-Ger\"ust
$\cal{E}^n$ angeheftet.

an kann daher einen CW-Komplex sukzessive aufbauen,
in dem man das $n$-Ger\"ust durch Anheften von $n$-cells
aus dem $n-1$-Ger\"ust konstruiert.
Beginnend with einem $0$-Ger\"ust (discrete set of points) entsteht der Komplex
durch Anheften of cells gr\"o\"sserer dimension.

The Begriff des \textit{Unterkomplexes} of a CW-Komplexes ergibt sich in
nat\"urlicher Weise.

{\bf Definition 15} (Unterkomplex).

Let $(\cal{E},\topX)$ a CW-Komplex,
$\cal{E}' \subset \cal{E}$ a subset of cells
and $\topX'$ the union ihrer points.
$(\cal{E}',\topX')$ is called a \textit{Unterkomplex} von $(\cal{E},\topX)$,
if $(\cal{E}',\topX')$ also a CW-Komplex ist.

Beliebige unions and beliebige intersections von Unterkomplexen sind
also Unterkomplexe.

\begin{examples}
\begin{itemize}
\item
Jedes $n$-Ger\"ust of a CW-Komplexes is a Unterkomplex.
\item
The union beliebig vieler $n$-cells $\epsilon \xeof \cal{E}$
with dem $(n-1)$-Ger\"ust $\cal{E}^{n-1}$ is a Unterkomplex.
\end{itemize}
\end{examples}

Unterkomplexe lassen sich wie folgt charakterisieren.

{\bf Satz 4}

Let $(\cal{E},\topX)$ a CW-Komplex.
Ein CW-Komplex $(\cal{E}',\topX')$ is genau then a Unterkomplex von
$(\cal{E},\topX)$,
if $\cal{E}' \subset \cal{E}$ and $\topX' \subset \topX$ in $\topX$
abgeschlossen ist.

Ein CW-Komplex is called endlich,
if $\cal{E}$ nur aus endlich vielen cells besteht.
A CW-complex is compact iff it is finite.

Jede Zelle $\epsilon$ liegt in einem endlichen Unterkomplex von $\cal{E}$.
Jede kompakte subset of a CW-Komplexes ist
in einem endlichen Unterkomplex enthalten.

\fancyhead[C] {\small{\textsc{VI. Singular Homology}}}
%------------------------------------------------------------------------------%
\part{Singular Homology}
%------------------------------------------------------------------------------%
A \textit{homology theory} is a topological-algebraical functor,
which assigns to each topological space a chain complex.
The diesem chain complex zugrundeliegende Ring $\riR$ wird als commutative
vorausgesetzt and owns a unit element.

The am meisten betrachtete homology theory is the \textit{singular} homology,
die for each topological space erkl\"art werden kann.
Im Mittelpunkt steht hierbei der Begriff der \textit{singular chain} and
ihre Randabbildungen.

\section{Singular Chains}
%------------------------------------------------------------------------------%
Ein $n$-dimensionale \textit{Standard-Simplex} im $\R^{n+1}$ is erkl\"art
als the set
\[
\Delta^n := \set{(x_0,\ldots,x_n) \xeof \R^{n+1}}
                { \, 0 \leq x_i \leq 1, \sum_{i=0}^n x_i \eq 1}.
\]
Speziell for $n=0$ sei $\Delta^0 := \{0,1\}$.

Setzt man a Koordinate $x_i \eq 0$,
so erh\"alt man a $(n-1)$-dimensionale Seite $\delta_i \Delta^n$
von $\Delta^n$.
Diese Seite kann man als Bild einer Einh\"angungs-mapping
$\delta^n_i : \Delta^{n-1} \to \Delta^n$ mit
\[
\delta^n_i(x_0,\ldots,x_{n-1}) \eq (x_0,\ldots,x_{i-1},0,x_i,\ldots,x_n)
\]
auffassen.
Diese mapping $\delta^n_i$ is called the $i$-te \textit{Seitenabbildung} von
$\Delta^{n-1}$ nach $\Delta^n$.

The Seitenabbildungen gen\"ugen gewissen Vertauschungsregeln.
For each $j < i$ gilt
$\delta^n_i \circ \delta^n_j \eq \delta^n_j \circ \delta^n_{i-1}$,
w\"ahrend for $j \geq i$ the Identit\"at
$\delta^n_i \circ \delta^n_j \eq \delta^n_{j+1} \circ \delta^n_i$ gilt.

it Hilfe der Standard-Simplices and ihre Seitenabbildungen kann man den
Begriff des \textit{singular simplices} in general topological spaces
einf\"uhren:

{\bf Definition 1} (singular $n$-simplex).

Ein {singular $n$-simplex} is a continuous mapping
$\sigma: \Delta^n \to \topX$ des $n$-dimensionalen
Standard-Simplices in a topological space $\topX$.

The Seitenabbildungen $\delta^n_i$ between Standard-Simplices induzieren die
Definition von Seitenabbildungen between singular Simplices:

{\bf Definition 2} (Seite of a singular $n$-Simplices).

Let $\topX$ a topological space,
$S_n(\topX)$ the set aller singular $n$-Simplizes in $\topX$
and $\sigma \xeof S_n(\topX)$.
Then $\partial_i \sigma := \sigma \circ \delta_i^n \xeof S_{n-1}(\topX)$
is called the \textit{$i$-te Seite} of $\sigma$.

The Seitenbildung of an element $ \sigma \xeof S_n(\topX)$ fulfills the identity
$(\partial_i \circ \partial_{i-1}) \sigma \eq 0$.

To apply methods from homological algebra,
dehnt man the set $S_n(\topX)$ to a $\riR$-module $\cal{S}_n(\topX;\riR)$ aus,
for den $S_n(\topX)$ a free basis darstellt.

{\bf Definition 3} (singular $n$-chain).

Let $\cal{S}_n(\topX;\riR)$ der freie $\riR$-module with basis $S_n(\topX)$.
Then are called the elements von $\cal{S}_n(\topX;\riR)$ \textit{singular $n$-chains}
in $\topX$ (with coefficients in $\riR$).

The singular $n$-chains are also formale $\riR$-Linearkombinationen von
singular Simplices.

The Seitenabbildungen $\partial_i$ between singular Simplices
gestatten \"uberdies the Definition von module-homomorphisms
between den modules $\cal{S}_n(\topX;\riR)$:
sei $\Rand{n}: \cal{S}_n(\topX;\riR) \to \cal{S}_{n-1}(\topX;\riR)$
die mapping mit
\[
\Rand{n}(\sigma) \, := \, \sum_{i=0}^n (-1)^i \partial_i \sigma.
\]
Jede mapping $\Rand{n}$,
welche der Komplex-Bedingung $\Rand{n-1} \circ \Rand{n} \eq 0$ gen\"ugt,
ist a module-homomorphism.
This implies
that the sequence of $\riR$-modules $\{\cal{S}_n(\topX;\riR)\}$ with the
module-homomorphisms $\Rand{n}$ forms a $\riR$-chain complex $K(\topX;\riR)$.
Each homomorphism $\Rand{i}$ is called a \textit{boundary operator}
of this chain complex.

The Untersuchung der homology-groups von $K(\topX;\riR)$ is naheliegend.

{\bf Definition 4} (singular homology-module).

Let $\topX$ a topological space.
Then is called $H_n(\topX;\riR)$ der $n$-te \textit{singular homology-module} von $\topX$.

it $H_{\*}(\topX,R) := \bigoplus_{n \geq 0} H_n(\topX;\riR)$ bezeichnet man die
\textit{totale singular homology} des topological space $\topX$.

\begin{example}
Let $\riR=\Z$ and $\topX \eq \{p\}$ a singleton.
Then gibt es for each $n \xeof \N$ genau a singular $n$-Simplex.
Es gilt also $\cal{S}_n(\topX;\Z) \cong \Z$.
For the Randoperatoren erh\"alt man $\Rand{i} \eq 0$ for $i$ ungerade and
$\Rand{i} \eq Id$ for $i$ gerade.
Hieraus folgt insbesondere $H_0(\{p\},\Z) \cong \Z$.
\end{example}

Bis auf wenige simple examples are the singular chain complexes of a
topological space un\"uberschaubare and kaum berechenbare objects.
an ben\"otigt daher further methods,
to gain insight into the structure of these chain complexes.

\section{The Singular Homology-Functor}
%------------------------------------------------------------------------------%
Seien im folgenden $\topX,\topY$ topological spaces and
$f:\topX \to \topY$ a continuous mapping.
For each singular Simplex $\sigma \xeof S_n(\topX)$ is the mapping
$f \circ \sigma$ a singular Simplex von $\topY$.
Let $f^{\*}:S_n(\topX) \to S_n(\topY)$ the zugeh\"orige mapping.
Then gilt for $\sigma \xeof S_n(\topX)$:
\[
\Rand{n+1} \circ f^{\*}(\sigma) \eq f^{\*} \circ \Rand{n+1} (\sigma),
\]
that is,
$f$ induziert a homomorphism $\phi_f$ between den
chain complexes $K(\topX;\riR)$ and $K(\topY;\riR)$.

Ferner induziert $f$ for each $n \xeof \Z$ a module-homomorphism
\[
f^{\*}_n : H_n(\topX;\riR)
\to H_n(\topY;\riR), \quad f^{\*}([\sigma]) := [f \circ \sigma]
\]
between den homology-modules von $\topX$ and $\topY$.

{\bf Definition 5} (induzierte chainsabbildung).

$f^{\*}_n$ is called the von $f$ \textit{induzierte chainsabbildung}.

The mapping $f^{\*}_n$ stellt daher for each $n \xeof \N$ einen
kovarianter functor von der category der topological spaces
in the category der $\riR$-modules dar dar.

\section{Relationships Between Homology and Homotopy}
%------------------------------------------------------------------------------%
Seien $\topX,\topY$ topological spaces and $f:\topX \to \topY$ stetig.
Ist $f$ a homotopy-equivalence,
so is der homomorphism $\phi_f$ between den chain complexes
$K(\topX;\riR)$ and $K(\topY;\riR)$ kettenhomotop.

The singular homology is with der homotopytheorie insofern vertr\"aglich,
als dass homotopie-\"aquivalente spaces zueinander isomorphe homology-modules
besitzen.
Hieraus folgt for the induzierten chainsabbildungen:

{\bf Satz 1}

Ist $f:\topX \to \topY$ a homotopy-equivalence
between the spaces $\topX$ and $\topY$,
so are for each $n \xeof \Z$ the homology-modules
$H_n(\topX;\riR)$ and $H_n(\topY;\riR)$ zueinander isomorph.

Zwischen der Fundamentalgruppe and der ersten homology-Gruppe eines
topological space $\topX$ besteht folgender Zusammenhang:
ist $\gamma$ a Schleife an $p \xeof \topX$, so is $\gamma$ a $1$-Zykel,
that is,
$\gamma \xeof S_1(\topX)$ and hence $[\gamma] \xeof H_1(\topX,\riR)$.
The dadurch induzierte mapping $h:\pi_1(\topX,p) \to H_1(\topX,\riR)$ is ein
group homomorphism with $Kern h \eq Kom(\pi_1(\topX,p))$.
Hieraus folgt der Satz von Hurewicz:

{\bf Satz 2}

Let $\topX$ a wegzusammenh\"angender topological space.
\begin{itemize}
\item
Ist $\pi_1(\topX)$ kommutativ,
so are $\pi_1(\topX)$ and $H_1(\topX,\Z)$ zueinander isomorph.
\item
Ist $\pi_1(\topX)$ nicht-kommutativ,
then are $\pi_1(\topX)^{ab}$ and $H_1(\topX,\Z)$ zueinander isomorphic groups.
\end{itemize}

\section{The Homology Groups of Some Geometrical Spaces}
%------------------------------------------------------------------------------%
Let $\riR$ be a commutative and unital ring.
Then gilt
\[
H_m(S^n;\riR) \cong \riR, \quad m \eq n > 0, \quad m \eq 0
\]
and $H_0(S^0;\riR) \cong \riR \otimes \riR$ and $H_m(S^n;\riR) \cong 0)$ sonst.

Explizit bedeutet dies $H_0(S^0) \cong R$ and $H_0(S^1) \cong H_1(S^1) \cong \riR$.
Speziell for $n \eq 2$ erh\"alt man $H_0(S^2) \cong \riR$,
$H_1(S^2) \cong 0$ and $H_2(S^2) \cong \riR$.

Hieraus folgt,
dass Sph\"aren unterschiedlicher Dimension weder homotopie-\"aquivalent
noch hom\"oomorph zueinander sind.

For the real projective space we have
\[
H_i(\R\,P^n,\Z) \cong \Z, \quad i=0, \, i=2j-1
\]
and
\[
H_i(\R\,P^n,\Z) \cong \Z_2, \quad i=2j-1, \, i < n.
\]
In all other cases $H_i(\R\,P^n,\Z) \cong 0$ holds.
This means
\[
H_0(\R\,P^1,\Z) \cong \Z \cong H_1((\R\,P^1,\Z) \cong 0,
\]
\[
H_0(\R\,P^2,\Z) \cong \Z, \quad  H_1((\R\,P^2,\Z) \cong \Z_2,
\quad H_2(\R P^2,\Z) \cong 0.
\]

For the komplex projective space we have
\[
H_i(\C\,P^n,\Z) \cong \R, \quad i=2*j, \, 0 \leq j \leq n
\]
and $H_i(\C\,P^n,\Z) \cong 0$ otherwise.

\section{Eilenberg-Zilber Theorem}
%------------------------------------------------------------------------------%
For topological spaces is der singular chain complex des Produktraumes
in a natural way homotopie\"aquivalent zum Tensorprodukt der singular
Komplexe der Faktoren.
Dies is der Inhalt des \textit{Satzes von Eilenberg-Zilber}:

{\bf Satz 3} (Satz von Eilenberg-Zilber).

Let $\topX,\topY$ be topological spaces.
Then there are up to chain homotopy uniquely defined homomorphisms
\[
\phi\: S_{\*}(\topX) \otimes S_{\*}(\topY) \to S_{\*} (\topX \times \topY),
\]
\[
\psi\: S_{\*}(\topX \times \topY) \to S_{\*}(\topX) \otimes S_{\*}(\topY).
\]

\section{Cohomology}
%------------------------------------------------------------------------------%
Cohomology is the zur homology \textit{duale} Art and Weise,
topological spaces rsp. pairs of spaces $\riR$-modules zuzuweisen.
The Vorteil der cohomology besteht darin,
dass in der totalen cohomology auf nat\"urliche Weise an inner
multiplicative operation erkl\"art werden kann
such that
the total cohomology sogar a Ringstruktur tr\"agt.

In topology there is an important reason,
neben singular chain complexes auch the zugehorigen cochain complexes
zu betrachten.
Diese werden durch Dualisieren,
that is,
durch Anwenden of the \textit{Hom}-functor
for each module des chain complex gewonnen.
The so definierte cohomology tr\"agt neben der modulestruktur zus\"atzlich
eine Ringstruktur.
The Information, welche the cohomology over a topological space
liefert, is daher reichhaltiger als the der homology.
Dies is der Hauptgrund, in der Topologie cohomologies zu betrachten.

%{\bf Definition 6} (singular cochain, Korand-Operator, cohomology-module).
\begin{definition}
Let $\riR$ be a commutative unital ring, $\modM$ be a $\riR$-module and
$\topX$ be a topological space.

\begin{itemize}
\item
$S^n(\topX;\modM) := Hom_{\riR}(\cal{S}(\topX;\riR),\modM)$ is called der module
of $n$-dimensional \textit{singular cochains} of $\topX$
with values in $\modM$.
\item
The Korand-operator $d^n : S^n(\topX;\modM) \mapsto S^{n+1}(\topX;\modM)$
is defined according to
\[
(d^n(\phi)(\sigma) := (-1)^n \, \phi(\Rand{n+1}(\sigma))
\]
for all $\phi \xeof S^n(\topX;\modM)$ and all $\sigma \xeof S_{n+1}(\topX;\riR)$
as well as their linear extensions.
\item
$H^n(\topX;\modM) := Z^n(\topX;\modM) / B^n(\topX;\modM) \eq H^n(S(\topX;\modM),d)$
is called the $n$-te singular cohomology-module of $\topX$ with values in $\modM$.
\end{itemize}
\end{definition}

\subparagraph{}
The sequence $S^n(\topX;\modM),d^n)$ forms a cochain complex
due to $d^n \circ d^{n+1}$.
Let $f\:\topX \xto \topY$ be a continuous mapping.
Then $f^{\*}$ is called the \textit{dual mapping} for $f_{\*}$.
Hence for each $n \xeof \N$ $H^n$ is a contravariant functor
from the category of pairs of topological spaces
to the category of $\riR$-modules.

\section{Homotopy, Homology and Cohomology}
%------------------------------------------------------------------------------%
If $\topX,\topY$ are topological spaces,
then the following relationship between homeomorphy, homotopy-equivalence
and the isomorphy of homology and cohomology complexes holds:
\[
\topX \cong \topY
\quad
\Longrightarrow
\quad
\topX \simeq \topY
\quad
\Longrightarrow
\quad
H^{\star}(\topX,\Z) \cong H^{\star}(\topY,\Z)
\quad
\Longrightarrow
\quad
H_{\star}(\topX,\Z) \cong H_{\star}(\topY,\Z).
\]
The reverse statements do not hold.
since it is possible to give counterexamples for them.

%\end{comment}
%------------------------------------------------------------------------------%
% close document                                                               %
%------------------------------------------------------------------------------%
\end{document}
